\documentclass[11pt,a4paper]{article}
\usepackage[T1]{fontenc}
\usepackage[utf8]{inputenc}
\usepackage{lmodern,amsmath,amssymb,amsthm,mathtools}
\usepackage[margin=26mm]{geometry}
\usepackage{microtype,booktabs,longtable,array,enumitem,listings,xcolor}
\usepackage[colorlinks=true,linkcolor=blue!45!black,citecolor=blue!45!black,urlcolor=blue!45!black]{hyperref}
\setlength{\emergencystretch}{3em}
\setlist{itemsep=3pt,topsep=5pt}
\newtheorem{theorem}{Theorem}[section]
\newtheorem{lemma}[theorem]{Lemma}
\newtheorem{proposition}[theorem]{Proposition}
\newtheorem{corollary}[theorem]{Corollary}
\theoremstyle{remark}\newtheorem{remark}[theorem]{Remark}
\newtheorem{example}[theorem]{Example}
\newcommand{\Z}{\mathbb Z}\newcommand{\Q}{\mathbb Q}
\newcommand{\R}{\mathbb R}\newcommand{\F}{\mathbb F}
\newcommand{\Tr}{\operatorname{Tr}}\newcommand{\Nm}{\operatorname{N}}
\DeclareMathOperator{\ord}{ord}
\lstset{basicstyle=\ttfamily\scriptsize,breaklines=true,columns=fullflexible,keepspaces=true,showstringspaces=false}
\hypersetup{pdftitle={Cubic Pisot Floors: Composite Windows, Sparse Exceptions, and Finite-Sieve Barriers},pdfauthor={Keio}}
\title{Cubic Pisot Floors:\\Composite Windows, Sparse Exceptions,\\and Finite-Sieve Barriers}
\author{Keio}
\date{21 October 2026\\\small Revised research manuscript: residual fields and finite-prime survival}
\begin{document}
\maketitle
\begin{abstract}
We investigate Saito's question on composite values of
$\lfloor\beta^{3^n}\rfloor$ for cubic Pisot numbers, using the
supplied constant-race manuscripts as prior work. A constructive
prime-divisor argument places every fixed window of shifted floors
simultaneously on a composite arithmetic progression for fourteen
coefficient classes modulo three. A separate finite-field theorem
gives exactly $p^{d-1}$ coefficient vectors at every shifted trace
value for odd degree $d$ and primes $p\equiv2\pmod3$. Combined with
geometric separation and the large sieve, it yields a uniform
$\exp(-c\sqrt{\log m})$ exceptional proportion in coefficient boxes,
for every $c<1/(e\sqrt{2\log3})$. We also strengthen a Pell-family
exclusion by three exact modular certificates, refine conditional
finite prime/composite patterns to prescribed exact local
multiplicities, and count balanced Pisot units in fixed congruence
cosets. A residual-field analysis excludes every nonzero rational power
of Mills' constant from broad Galois and cyclotomic fields. Conversely,
in every totally real cubic field admitting a three-adic embedding,
we construct infinitely many Saito-balanced Pisot numbers avoiding
any prescribed finite set of prime divisors at every odd exponent.
Previously established resolutions and conditional prefix
results are identified explicitly. Neither the full cubic Pisot
problem, irrationality of the least square-clock constant, nor
unconditional transcendence of Mills' constant is proved.
\end{abstract}


\section{The problem, the source boundary, and the new results}
Let
\[
 A=\min\{x>1:\lfloor x^{3^k}\rfloor\text{ is prime for every }k\ge1\}
\]
be Mills' constant. Saito proved its irrationality in
\cite{SaitoIrrational}. His later reduction \cite[Theorem~1.7]{SaitoVariants} states that if $A$ is
algebraic, some $\beta=A^{3^m}$ is a totally real cubic Pisot number
whose two small conjugates satisfy
\begin{equation}\label{eq:saitobounds}
 |\beta_3|<-\beta_2\le
 \min\{|\beta_3|^{17/23},\beta^{-17/40}\}.
\end{equation}
The trace is then the floor at all sufficiently large cubic-power
levels. Excluding every such algebraic candidate would prove
unconditional transcendence of $A$. Saito's Problem~1.1 in
\cite{SaitoIntervals} asks the stronger question: does every cubic
Pisot number have infinitely many composite cubic-power floors?

The word ``question'' matters here. The source papers contain
explicit problems and alternatives, and an answer to a modified
sequence does not automatically settle the original least constant.
All source theorem numbers below refer to the specified arXiv
versions, since journal numbering may differ.

\subsection{What is already resolved and what remains}
\begin{center}\small
\begin{tabular}{p{.25\textwidth}p{.29\textwidth}p{.35\textwidth}}
\toprule
Source and item & Mathematical target & Status used in this paper\\
\midrule
\cite{SaitoIrrational}, Q1.6 & Irrationality of the least constant with exponent $2^k$ & Not proved here.\\
\cite{SaitoIrrational}, Q1.7 & The cubic Pisot power alternative for Mills' constant & Not excluded in full.\\
\cite{SaitoIrrational}, Q5.3 & RH implies Mills transcendence & Already answered by Saito; his later Theorem~1.8 uses the Density Hypothesis.\\
\cite{SaitoIntervals}, P1.1 & All cubic Pisot floors at $3^n$ & Restricted closures and quantitative exceptions below; full problem remains.\\
\cite{SaitoIntervals}, P1.7--1.8 & Nonreversible clocks and shifted dyadic Fibonacci values & Already treated in the supplied older manuscripts [008] and [007].\\
\cite{SaitoDimension}, Q1.4--1.5 & Unbounded multiplier and Wright-set dimensions & Already treated in supplied older [031], Theorems~14.3 and~15.1.\\
\bottomrule
\end{tabular}
\end{center}
The last two rows report the content and audited proofs of the
user-supplied manuscripts, not a claim of journal acceptance or
historical priority. They are not counted again as new resolutions.
The RH implication is a published external result of Saito, not a
consequence newly obtained from the present work.

\subsection{Contribution relative to the supplied manuscripts}
The new attachment inventory contains 244 TeX manuscripts, including
38 in the Mills--Pisot topic. We focused on their arithmetic proofs
and on plausible transfers from the other topics; this is not a
claim to have independently re-proved all 244 manuscripts.
The core additions and their limitations are as follows.
\begin{enumerate}
\item Theorem~\ref{thm:windows} synchronizes finitely many shifts on
one progression. The old single-shift positive-density statements
alone do not supply this synchronization.
\item Theorem~\ref{thm:shiftedsieve} extends the unshifted cubic
coefficient sieve of round S37 to every shift and every odd degree,
uniformly in the shift, with an improved explicit range of constants.
The older [031] already contains qualitative density-one results;
the quantitative rate and the stated extensions are the additions.
\item Theorem~\ref{thm:pell-new} adds three exact witnesses to the
Pell cover of rounds S35--S36. It removes two further valuation
layers, but leaves an infinite suborbit to be analyzed.
\item Theorem~\ref{fac:valuationpatterns} refines the finite patterns
of S28 by specifying exact selected prime valuations and their
periodic recurrence. Prime-designated positions still require
Schinzel's hypothesis H. The all-index irreducibility of S27 is
reproduced as prior input, not counted as new.
\item Theorem~\ref{transfer:thm:ray} quantifies the persistence of
Saito-balanced units in a fixed congruence coset, explaining why
finite local exclusions and the real size conditions can coexist.
\item Theorems~\ref{res:thm:ambient} and~\ref{res:thm:galois} turn
an inherited cubic local obstruction into exclusions for all rational
powers of Mills' constant in larger number fields. Theorem~
\ref{rescyc:thm:norm-one-mod-three} records the existing cyclic
norm restriction in a positive-density formulation for both floor branches.
\item Theorem~\ref{res:thm:equivalence} removes the seed assumption
from finite-prime survival by allowing a fixed nonzero integral norm.
It applies to every remaining totally real cubic field and every
prescribed finite prime set. Corollary~\ref{res:cor:normalized}
places such survivors in any real interval after normalization.
\end{enumerate}
The manuscript contains mathematical proofs and exact finite
certificates; external priority remains unestablished. In particular,
none of the remaining numbered global questions in the first four
rows is claimed as a new complete resolution.

\subsection{Conventions}
A composite integer is positive and larger than one, with a proper
prime divisor. Finite initial exceptions are harmless only after
growth has been proved. Lower density always refers to the index
$n$, not to the enormous integer value of the trace. A Pisot number
is a real algebraic integer greater than one with every other
conjugate strictly inside the unit circle. Traces count all conjugates
with multiplicity when polynomial specializations are used.


\section{Simultaneous composite windows from a three-adic limit}\label{sec:windows}
For a monic integral cubic $X^3-uX^2+vX-N$, let $U_n$ be its
$3^n$-th root power sum, let $V_n$ be the corresponding second
symmetric function, and let $E_n=N^{3^n}$. Newton's identities give
\begin{equation}\label{eq:cubicmap}
 (U_{n+1},V_{n+1},E_{n+1})=
 (U_n^3-3U_nV_n+3E_n,\,
 V_n^3-3E_nU_nV_n+3E_n^2,\,E_n^3).
\end{equation}
The norm coordinate must be advanced, except when $N=0,1,-1$.

\begin{lemma}[Three-adic contraction]\label{lem:contraction}
The sequence $U_n$ has a limit $\lambda\in\Z_3$, with
$\lambda-U_n\in3^{n+1}\Z_3$. This limit depends only on
$(u,v,N)\pmod3$.
\end{lemma}
\begin{proof}
Let $F$ denote the polynomial map in \eqref{eq:cubicmap}. It is the
identity modulo three. If $\mathbf x\equiv\mathbf y\pmod{3^r}$,
$r\ge1$, expansion gives
$F(\mathbf x)\equiv F(\mathbf y)\pmod{3^{r+1}}$; every first
derivative has coefficients divisible by three, and higher terms
have the requisite divisibility. Consecutive iterates therefore
differ by a vector divisible by $3^{n+1}$ at step $n$. The same
argument for two initial points congruent modulo three proves
equality of their limits.
\end{proof}

\begin{theorem}[A common progression for finitely many shifts]
\label{thm:windows}
Suppose $U_n\to+\infty$ in the real topology. Let
$\mathcal H\subset\Z$ be finite and assume
\[
 \lambda+h\notin\{1,-1\}\qquad(h\in\mathcal H).
\]
There are integers $n_0,D\ge1$ such that all
$U_{n_0+jD}+h$, $h\in\mathcal H$, are composite for every
sufficiently large $j$.
\end{theorem}
\begin{proof}
Choose a common large $n_0$ so that $Q_h=U_{n_0}+h>1$ for all
shifts. When $3\nmid\lambda+h$, also arrange
\[
 v_3(Q_h-1)<n_0,\qquad v_3(Q_h+1)<n_0.
\]
The two limiting valuations are finite by hypothesis, so
Lemma~\ref{lem:contraction} permits this choice.
If $3\mid\lambda+h$, the prime three divides $U_n+h$ at every
level and supplies a recurring divisor.

For any other shift factor $Q_h$ into primes. If every prime factor
$q$ satisfied $v_3(q^2-1)\ge n_0$, then each such $q$ would be
$1$ or $-1$ modulo $3^{n_0}$. Their product, including
multiplicities, would make $Q_h\equiv\pm1\pmod{3^{n_0}}$,
a contradiction. Hence there is a prime $q_h\mid Q_h$ with
\[
 v_3(q_h^2-1)\le n_0-1.
\]
This prime is not three. Every nonzero root of the cubic modulo
$q_h$ has order dividing
$M_h=\operatorname{lcm}(q_h-1,q_h^2-1,q_h^3-1)$.
Elementary lifting of valuations gives
\[
 v_3(M_h)\le v_3(q_h^2-1)+1\le n_0.
\]
After powering to $3^{n_0}$, all remaining root orders are prime
to three. A period $D_h$ for cubing on these groups therefore gives
$q_h\mid U_{n_0+jD_h}+h$ for every $j\ge0$. Zero roots and repeated
roots cause no difficulty; the power sum counts multiplicities.

Take $D$ to be the least common multiple of the finitely many
$D_h$, with $D=1$ when there are none. All chosen prime divisors recur on the same progression.
Real growth makes them proper divisors for all sufficiently large
$j$. Unlike a prime-seed argument, this construction works even
when every initial value selected is already composite.
\end{proof}

\begin{corollary}[Windows around cubic Pisot floors]\label{cor:pisotwindows}
Let $\beta$ be a cubic Pisot number and let $\lambda$ be its
three-adic trace limit. If $\lambda$ is not an ordinary integer,
then for every $H\ge0$ there is an arithmetic progression of indices
on which, eventually, all
\[
 \lfloor\beta^{3^n}\rfloor+h,\qquad |h|\le H,
\]
are composite simultaneously. Nonreal conjugates are permitted.
\end{corollary}
\begin{proof}
Apply Theorem~\ref{thm:windows} to all integer shifts in
$[-H-1,H]$. None can give $\lambda+h=\pm1$. The sum of the two
small conjugate powers tends to zero, so the floor is eventually
either $U_n$ or $U_n-1$. The stated floor window lies inside the
chosen trace window, regardless of changes in the sign of the
small-conjugate sum.
\end{proof}

Separate positive-density statements for different shifts cannot
be intersected without an argument. The common initial index and
the prime-divisor construction above provide that missing step.
The single-shift order-budget mechanism was already present in the
older manuscript \cite{Keio009}; the simultaneous progression is
the strengthening here.

\begin{proposition}[Exact residue-class scope]\label{prop:fourteen}
The limit $\lambda$ is an ordinary integer precisely when the
residual cubic over $\F_3$ either splits completely or is
$(X-r)(X^2+1)$ with $r\in\F_3$. Thus
Corollary~\ref{cor:pisotwindows} applies to fourteen of the
twenty-seven coefficient triples modulo three.
\end{proposition}
\begin{proof}
For each nonzero residual root, take its Teichm\"uller lift, the
unique prime-to-three root of unity with that residue. In a
three-adic splitting field an integral root of positive valuation
vanishes under repeated cubing; a unit root's principal-unit part
tends to one. Frobenius permutes the multiset of residual roots,
so $\lambda$ is the sum of the Teichm\"uller lifts, with
multiplicity.

A linear root contributes $0,1$, or $-1$. The irreducible quadratic
$X^2+1$ contributes two fourth roots of unity with sum zero. The
other two irreducible quadratics are $X^2+X-1$ and $X^2-X-1$.
Their lifted pairs have sum with square $-2$, so adding one linear
contribution cannot give an integer.

For an irreducible cubic, the nonzero roots have order thirteen or
twenty-six. The lifted trace is
$\pm(\zeta^a+\zeta^{3a}+\zeta^{9a})$, where $\zeta$ is a
primitive thirteenth root and $a\ne0\pmod{13}$. This is not
rational: otherwise a polynomial of degree at most twelve with
at most four nonzero coefficients would vanish at $\zeta$, in
contradiction to its minimal polynomial
$1+X+\cdots+X^{12}$. There are ten completely split monic cubics
(multisets of three roots from a three-element set), three of the
second integer-limit type, six with the other quadratic factors,
and eight irreducible cubics. These exhaust the twenty-seven cases.
\end{proof}

\begin{center}
\begin{tabular}{c l}
\toprule
Integer limit $\lambda$ & $(u,v,N)\pmod3$\\
\midrule
$-3$ & $(0,0,2)$\\
$-2$ & $(1,1,0)$\\
$-1$ & $(2,0,0),\ (2,1,2),\ (2,2,1)$\\
$0$ & $(0,0,0),\ (0,1,0),\ (0,2,0)$\\
$1$ & $(1,0,0),\ (1,1,1),\ (1,2,2)$\\
$2$ & $(2,1,0)$\\
$3$ & $(0,0,1)$\\
\bottomrule
\end{tabular}
\end{center}
The complement of this table is the fourteen-class scope. This is
a count of residue classes, not a density assertion about Pisot
numbers under an unspecified ordering. For example, the plastic
number, the real root of $X^3-X-1$, has irreducible reduction modulo
three and satisfies the simultaneous-window conclusion.

For the unshifted trace problem only the six classes with
$\lambda=\pm1$ can escape Theorem~\ref{thm:windows}. These include
the surviving Mills alternatives. For the predecessor trace
$U_n-1$, only the four classes with $\lambda=0$ or $2$ can escape;
all four have $N\equiv0\pmod3$. In particular, whenever
$3\nmid N$ and the floor is eventually $U_n-1$, its composite
indices have positive lower density. This last single-shift
consequence was already available in the earlier local obstruction.

% This section follows windows_section.tex and uses its labels and notation.
\section{Field criteria and a uniform density bound}
\label{local:sec:fields}
The preceding simultaneous-return theorem has a useful field-theoretic
consequence. We first record an elementary bound that also handles
changing signs of nonreal conjugates.

\begin{proposition}[A parity and trace criterion]
\label{local:prop:half}
Let $\beta$ be a cubic Pisot number of odd norm. If
$\operatorname{Tr}(\beta)\equiv0\pmod3$, then
\[
 \liminf_{M\to\infty}\frac{1}{M}
 \#\{0\le n<M:\lfloor\beta^{3^n}\rfloor
                  \text{ is composite}\}\ \ge\ \frac12.
\]
The two small conjugates may be nonreal.
\end{proposition}
\begin{proof}
Let $U_n,V_n,E_n$ be the coefficient sequences in
\eqref{eq:cubicmap}. Since the norm is odd, $E_n=1$ modulo two.
The map on $(U_n,V_n)\pmod2$ is
\[
 (a,b)\longmapsto(a+ab+1,\ b+ab+1).
\]
Its four states satisfy
\[
 00\longmapsto11,\qquad11\longmapsto11,
 \qquad01\longleftrightarrow10.
\]
Consequently $U_n$ is odd either at every sufficiently large index
or at exactly one parity of the indices. At least half of all
indices therefore have $U_n$ odd. At the same time,
\eqref{eq:cubicmap} modulo three gives $U_n\equiv U_0\equiv0\pmod3$
for every $n$.

The sum of the two small conjugate powers is real and tends to zero.
Thus, for every sufficiently large $n$,
$\lfloor\beta^{3^n}\rfloor$ belongs to $\{U_n,U_n-1\}$.
When $U_n$ is odd, the first possibility is divisible by three and
the second is divisible by two. Growth makes these divisors proper.
This proves the density bound without any assumption on the sign
pattern of the small-conjugate sum.
\end{proof}

\begin{theorem}[Cubic fields without a three-adic embedding]
\label{local:thm:noembedding}
Let $K$ be a cubic number field admitting no field embedding into
$\mathbb Q_3$. Let $\beta\in\mathcal O_K$ be a cubic Pisot number
with $3\nmid N_{K/\mathbb Q}(\beta)$. Then the composite indices of
$\lfloor\beta^{3^n}\rfloor$ contain, up to finitely many terms,
an arithmetic progression. In particular they have positive lower
density.

More precisely, write $f$ for the minimal polynomial of $\beta$.
If $f\bmod3$ is irreducible, every finite window about these floors
is simultaneously composite on a progression, as in
Corollary~\ref{cor:pisotwindows}. Otherwise $f\bmod3$ is a cube
of a nonzero linear factor. In this second case, if the norm is
also odd, the lower density is at least $1/2$.
\end{theorem}
\begin{proof}
A simple root of $f\bmod3$ would lift by Hensel's lemma to a root
in $\mathbb Q_3$. Since $K=\mathbb Q(\beta)$, that root would define
an embedding $K\hookrightarrow\mathbb Q_3$, contrary to hypothesis.
A monic cubic over $\mathbb F_3$ without a simple linear factor
is either irreducible or $(X-a)^3$. The norm assumption makes
$a\ne0$ in the latter case.

In the irreducible case Proposition~\ref{prop:fourteen} gives a
noninteger three-adic trace limit, and
Corollary~\ref{cor:pisotwindows} proves the asserted window statement.
In the remaining case the three Teichm\"uller lifts of the residual
roots are all $1$ or all $-1$. Thus the trace limit is $\lambda=3$
or $\lambda=-3$, respectively. Neither $\lambda$ nor $\lambda-1$
equals $1$ or $-1$. Apply Theorem~\ref{thm:windows} to the two
shifts $\{0,-1\}$. Both $U_n$ and $U_n-1$ are composite on one
eventual progression; the Pisot estimate again places the floor
among these two integers. Finally $(X-a)^3$ has zero quadratic
coefficient modulo three, so the odd-norm case satisfies
Proposition~\ref{local:prop:half}.
\end{proof}

\begin{corollary}[Totally ramified cubic fields]
\label{local:cor:ramified}
If three is totally ramified in a cubic number field $K$, then
every cubic Pisot unit of $K$ has composite floors at a set of
indices of lower density at least $1/2$.
\end{corollary}
\begin{proof}
The algebra $K\otimes_{\mathbb Q}\mathbb Q_3$ is a field of
degree three, so $K$ has no embedding into $\mathbb Q_3$.
For any $\beta\in\mathcal O_K$, its three local conjugates have
the same residue in the residue field $\mathbb F_3$. Its minimal
polynomial therefore reduces to $(X-a)^3$ modulo three. A unit
has nonzero residue and norm $\pm1$. The last assertion of
Theorem~\ref{local:thm:noembedding} applies.
\end{proof}

For example, $K=\mathbb Q(r)$ with $r^3-3r-1=0$ is a totally
real cyclic cubic field: the polynomial is irreducible, has
discriminant $81$, and its translate
$X^3+3X^2-3$ is Eisenstein at three. Hence the corollary covers
every Pisot unit in this field. It also covers every Pisot unit in
$\mathbb Q(\sqrt[3]{2})$, since
$(X-1)^3-2=X^3-3X^2+3X-3$ is Eisenstein at three.
An explicit example in the latter field is
$\beta=1+\sqrt[3]{2}+\sqrt[3]{4}$, whose minimal polynomial is
$X^3-3X^2-3X-1$. It has norm one and a unique real root greater
than one; its conjugate pair has modulus $\beta^{-1/2}<1$.
Thus the density bound applies even though the floor correction
need not have an eventually constant sign.

\paragraph{Completion of the norm-positive case in the field of
discriminant $1524$.}
Let $r$ be the largest root of $X^3-7X^2-X+1$ and let
$K=\mathbb Q(r)$. Showa round 20 proves that every norm-one unit
$\beta\in K$ with $\operatorname{Tr}(\beta^{3^n})\to+\infty$
has a positive lower density of composite trace terms. This is a
prior result of the supplied archive, not a new trace theorem here.
It also yields the following precise floor statement.

\begin{corollary}\label{local:cor:1524}
Every cubic Pisot unit of norm one in the field of discriminant
$1524$ has a positive lower density of composite floors
$\lfloor\beta^{3^n}\rfloor$.
\end{corollary}
\begin{proof}
The field is totally real. Since $\beta>1$ and its norm is one,
its two small conjugates have the same sign. If both are negative,
the floor equals the trace eventually, and the cited round-20
theorem applies. If both are positive, the floor equals $U_n-1$
eventually. The four-state parity calculation in the proof of
Proposition~\ref{local:prop:half} shows that $U_n$ is odd at a
set of indices of density at least $1/2$. At those sufficiently
large indices $U_n-1$ is a positive even composite integer.
\end{proof}

\paragraph{Scope and comparison with the supplied manuscripts.}
The old manuscript \cite{Keio009} already proves infinitely many
composite floors for every nonreal cubic Pisot number, and gives
positive-density results for fixed real shifts. Its projective
return theorem already includes the cubic-residue norm refinement;
that refinement is not new here. The absence of a three-adic
embedding was already a necessary obstruction for a hypothetical
prime trace tail in \cite{Keio031}. The new use here is the
simultaneous treatment of the two possible floor corrections,
together with the explicit density-$1/2$ criterion. In particular
the statements above do not depend on an unjustified intersection
of separate positive-density sets.

The discriminant-$1524$ corollary is a completion from traces to
all norm-one Pisot floors using a prior field-specific theorem.
It makes no assertion about norm-minus-one units in that field:
round 17 exhibits such units that survive several local tests.
Nor do these criteria exclude every totally real cubic Pisot
candidate in Saito's Problem~1.1 or establish Mills transcendence.

\section{Shifted traces: exact local counts and sparse exceptions}\label{sec:density}
This section extends the unshifted cubic sieve of submitted round
S37. The extensions are uniformity in the integer shift, both floor
branches in totally real cubic fields, and arbitrary odd degree. We
include the proof, since changing a trace to a shifted trace is not
automatically legitimate in a fixed-norm argument.

For an odd integer $d\ge3$, write
\[
 f_{\mathbf a}(X)=X^d-a_1X^{d-1}+a_2X^{d-2}-\cdots-a_d,
 \qquad P_n(\mathbf a)=\sum_{j=1}^d\alpha_j^{3^n}.
\]
Here the roots are counted with multiplicity, and Newton's identities
define $P_n\in\Z[a_1,\ldots,a_d]$ even at reducible specializations.
Its total degree is $3^n$, and its unique term of that degree is
$a_1^{3^n}$. For a sequence of integers, a \emph{recurring-divisor
certificate} is a prime dividing all terms on an infinite arithmetic
progression of indices. If the terms tend to positive infinity, such
a certificate gives positive lower density of composite terms.

\begin{theorem}[Uniform shifted sieve]\label{thm:shiftedsieve}
Put $c_0=1/(e\sqrt{2\log3})$. Fix an odd $d\ge3$ and $0<c<c_0$.
Let $\mathcal B\subset\Z^d$ be any box with side lengths
$X_1,\ldots,X_d\ge3$, and put $m=\min X_i$ and $V=\prod X_i$.
For every integer $h$, all but
\begin{equation}\label{eq:exceptionbound}
 O_{d,c}\bigl(V\exp(-c\sqrt{\log m})\bigr)
\end{equation}
coefficient vectors in $\mathcal B$ have a recurring-divisor
certificate for $P_n(\mathbf a)+h$. The constant is independent of
$h$ and of the location of the box.

There is a constant $C_{d,c}>0$ such that the certificate can be
chosen with prime $p\le\exp(C_{d,c}\sqrt{\log m})$ and progression
period at most $\exp(dC_{d,c}\sqrt{\log m})$. Consequently, whenever
$P_n(\mathbf a)\to+\infty$, the lower density of composite shifted
terms is positive, and is at least the reciprocal of that period.
\end{theorem}

The theorem gives an upper bound on a set of possible exceptions; it
does not prove that this set is empty. The assertion is about the
full coefficient space. A thin fixed field, fixed norm, or selected
one-parameter family requires a separate argument.

\begin{lemma}[Exact shifted local count]\label{lem:shiftedcount}
For every prime $p\equiv2\pmod3$, $n\ge1$, and $t\in\F_p$,
\[
 \#\{\mathbf a\in\F_p^d:P_n(\mathbf a)=t\}=p^{d-1}.
\]
Repeated roots and zero constant terms are allowed.
\end{lemma}
\begin{proof}
The coefficient vectors are Frobenius-stable orbits of ordered root
tuples under $S_d$. Twisted orbit counting gives the number in question
as $d!^{-1}\sum_{\sigma\in S_d}\#\operatorname{Fix}(F\sigma)$,
where the fixed tuples also satisfy the power-sum condition. This is
an identity of integer counts, including orbits with stabilizers.
All tuples under consideration lie in a common finite extension, so
ordinary finite-set counting applies.

A cycle of length $j$ in $\sigma$ corresponds to an element
$b\in\F_{p^j}$ and contributes
$\operatorname{Tr}_{\F_{p^j}/\F_p}(b^{3^n})$ to the power sum.
Every permutation of an odd number of letters has an odd cycle.
For such a cycle $j$, $3\nmid p^j-1$, so $b\mapsto b^{3^n}$ is a
bijection of $\F_{p^j}$. The field trace is surjective, with fibers
of size $p^{j-1}$. Fixing all other cycles leaves exactly that many
choices for this one, giving $p^{d-1}$ fixed tuples for each
permutation. Averaging proves the claim. Surjectivity of field trace
does not require $p\nmid j$.
\end{proof}

\begin{lemma}[Geometric separation]\label{lem:shiftedgeometry}
Over an algebraically closed field of characteristic different from
three, $P_n-t$ is irreducible for every $n\ge1$ and every $t$.
For distinct indices $i,j$ and any $s,t\in\F_p$, $p\ne3$,
\[
 \#\{P_i=s,\ P_j=t\}\le 3^{i+j}p^{d-2}.
\]
\end{lemma}
\begin{proof}
Under the finite, surjective Vieta map the first polynomial pulls
back to $Z_1^r+\cdots+Z_d^r-t$, $r=3^n$. Regard this as a polynomial
in $Z_d$ over $k[Z_1,\ldots,Z_{d-1}]$. The nonconstant polynomial
$Z_1^r+\cdots+Z_{d-1}^r-t$ has a prime factor with multiplicity one:
it is squarefree, since any repeated factor would divide all its
partial derivatives, which are multiples of the relatively prime
monomials $Z_1^{r-1},\ldots,Z_{d-1}^{r-1}$. Eisenstein at such a
factor proves irreducibility of the pullback, and hence of $P_n-t$.

Different indices have different degrees, so the two hypersurfaces
have no common component. Their intersection has pure dimension
$d-2$ and degree at most $3^{i+j}$, by B\'ezout. An affine variety of
dimension $r$ and degree $D$ has at most $Dp^r$ rational points.
One proof is induction on dimension: on each irreducible component
choose a nonconstant coordinate, slice at its $p$ rational values,
and apply the degree bound for hyperplane sections. Constant
coordinates are discarded. This argument also handles components
defined only over an algebraic closure, since it bounds their
$\F_p$-points individually. These facts give the displayed bound.
\end{proof}

\begin{lemma}[A uniform return budget]\label{lem:returnbudget}
Fix $B\ge1$ and put $b_d=\lfloor\log_3(d/2)\rfloor$.
Suppose $p\equiv2\pmod3$ and $v_3(p+1)\le B$. For every coefficient
vector modulo $p$, its power-sum sequence is periodic from index
$B+b_d$, with period less than $p^d$.
\end{lemma}
\begin{proof}
For an irreducible root factor of degree $j$, each nonzero root has
order dividing $p^j-1$. For odd $j$ this order is prime to three;
for even $j=2k$, the lifting-the-exponent identity gives
\[
 v_3(p^{2k}-1)=v_3(p+1)+v_3(k)\le B+b_d.
\]
After that many cubings all nonzero root powers have orders prime
to three. Cubing is then a permutation of their finite groups.
Zero roots stay zero and multiplicities do not change this argument.
The least common multiple $M$ of the remaining root orders is at
most the product of $p^j-1$ over distinct nonzero irreducible factors,
which is less than $p^d$. A common period is the order of $3$ modulo
$M$, with period one when $M=1$.
\end{proof}

\begin{proof}[Proof of Theorem~\ref{thm:shiftedsieve}]
The eligible primes of Lemma~\ref{lem:returnbudget} form fixed
reduced residue classes modulo $3^{B+1}$ of density
\[
 \delta=\delta_B=\frac{1-3^{-B}}2.
\]
Let $k_0=B+b_d$ and use the $L$ consecutive levels
$k_0,\ldots,k_0+L-1$. At an eligible prime let $\Omega_p$ be the
union of the loci $P_n+h=0$ at these levels, and put
$\omega(p)=|\Omega_p|$. Exact counts and the first two
inclusion--exclusion inequalities give
\[
 \frac Lp-\frac{D_L}{p^2}\le\frac{\omega(p)}{p^d}\le\frac Lp,
 \qquad
 D_L=\sum_{k_0\le i<j<k_0+L}3^{i+j}<9^{k_0+L}.
\]
For $p>y:=9^{k_0+L+1}$ we therefore have
\begin{equation}\label{eq:uniformforbidden}
 \frac{L-1}{p}\le\frac{\omega(p)}{p^d}\le\frac Lp.
\end{equation}
These estimates are independent of the shift, even when it varies
with the box.

We use the classical multidimensional large sieve in the form
\begin{equation}\label{eq:largesieve}
 |S|\ll_d\frac{(m+Q^2)^d}{H(Q)},\qquad
 H(Q)=\sum_{\substack{q\le Q\\q\ {\text{squarefree}}}}
 \prod_{p\mid q}\frac{\omega(p)}{p^d-\omega(p)},
\end{equation}
for a set $S$ in a cube of side $m$ avoiding the forbidden classes.
For completeness, its normalization can be seen directly. On
$(\Z/p\Z)^d$ take a mean-zero function equal to $1$ on the allowed
classes and to $-(p^d-\omega(p))/\omega(p)$ on the forbidden classes.
Its normalized squared norm is $1/g(p)$, where
$g(p)=\omega(p)/(p^d-\omega(p))$. The Chinese-remainder product,
Parseval, and Cauchy--Schwarz imply
\[
 g(q)|S|^2\le
 \sum_{\substack{\mathbf a\bmod q\\
 \gcd(a_1,\ldots,a_d,q)=1}}
 \left|\sum_{\mathbf x\in S}e^{2\pi i\mathbf a\cdot\mathbf x/q}\right|^2.
\]
Primitive frequencies with denominators at most $Q$ are
$Q^{-2}$-separated. The analytic large sieve bounds their combined
energy by $O_d((m+Q^2)^d)|S|$, proving
\eqref{eq:largesieve}. This is a classical input, not a new sieve
inequality; see \cite{Kowalski}.

Set $Q=\lfloor\sqrt m\rfloor$, $s=\sqrt{\log Q}$,
$L=\lfloor a s\rfloor$, and $z=\exp(bs)$, with fixed positive
$a,b$ such that
\[
 b>a\log9,\qquad ab\delta<1.
\]
Only eligible primes in $(y,z]$ are used. The prime number theorem
in these fixed residue classes and partial summation give
\[
 \sum_{y<p\le z}^{*}\frac1p
 =\delta\log\frac{b}{a\log9}+o(1),\qquad
 \sum_{p\le z}^{*}\frac{\log p}{p}=(\delta+o(1))\log z.
\]
The star restricts to eligible primes. The moving endpoints cause
no uniform-modulus issue: $B$ is fixed first and both endpoints
tend to infinity.

Expand the finite Euler product $W=\prod_{y<p\le z}^{*}(1+g(p))$
with its nonnegative squarefree weights. Equation
\eqref{eq:uniformforbidden} implies
\[
 \log W\ge
 \left(a\delta\log\frac b{a\log9}+o(1)\right)s.
\]
The weighted mean of $\log q$ in this expansion is exactly
\[
 \sum_{y<p\le z}^{*}\frac{g(p)}{1+g(p)}\log p
 \le (ab\delta+o(1))\log Q.
\]
Since $ab\delta<1$, Markov's inequality leaves a fixed positive
fraction of the weight on $q\le Q$. Thus $H(Q)\gg W$.
Take
\[
 a=\frac1{e\sqrt{\delta\log9}},\qquad
 b=(1-\varepsilon)e\sqrt{\frac{\log9}{\delta}}
\]
with $\varepsilon>0$ sufficiently small. As $\varepsilon\downarrow0$,
the coefficient in the lower bound for $\log H(Q)$, when expressed
in units of $\sqrt{\log m}$, approaches
$e^{-1}\sqrt{\delta/\log3}$. Finally choose a fixed large $B$ so
that this exceeds the prescribed $c<c_0$. We obtain
$H(Q)\gg_{d,c}\exp(c\sqrt{\log m})$.

The cube sieve proves \eqref{eq:exceptionbound} for a cube. A
rectangular box is covered by at most $2^dV/m^d$ such cubes. A point
not in the sifted set has a zero at one of the chosen levels modulo
some $p\le z$. Lemma~\ref{lem:returnbudget} supplies the claimed
certificate and period bound. Bounded values of $m$ are absorbed
in the constants.
\end{proof}

\begin{corollary}[Both totally real floor branches]\label{cor:floorbox}
Among coefficient triples in any box of minimum side $m$, the
totally real cubic Pisot polynomials for which
$\lfloor\beta^{3^n}\rfloor$ fails to be composite on a set of
positive lower density number at most
\[
 O_c\bigl(X_1X_2X_3e^{-c\sqrt{\log m}}\bigr),\qquad c<c_0.
\]
\end{corollary}
\begin{proof}
If the two small real conjugates are $\gamma,\eta$, their odd-power
sum has a constant nonzero sign for all sufficiently large powers.
The only potential exception is $\gamma=-\eta$, which would make
$\beta$ the integer trace, contradicting cubic irreducibility.
Thus the floor is eventually $P_n$ or $P_n-1$. Apply the theorem
to both shifts and take a union bound. The traces tend to infinity.
\end{proof}

\begin{corollary}[Many shifts]\label{cor:manyshifts}
For any finite set $\mathcal H\subset\Z$, the exceptional count for
having a certificate for every $h\in\mathcal H$ is at most
$O_{d,c}(|\mathcal H|V e^{-c\sqrt{\log m}})$.
In particular, for $0<\eta<c<c_0$, one may take
$|h|\le\exp(\eta\sqrt{\log m})$ and retain a saving
$e^{-(c-\eta)\sqrt{\log m}}$.
The progressions for different shifts in this corollary need not
coincide. A genuinely simultaneous progression is proved separately
in Section~\ref{sec:windows}.
\end{corollary}

\section{Two further valuation layers on the twice-square Pell locus}
\label{sec:pell-new}
Consider
\[
 g_u(X)=X^3-uX^2-X+1,\qquad
 4u^2+5u+23=2w^2,\quad u\ge2,\quad w>0.
\]
Let $P_n(u)$ be the sum of the $3^n$-th powers of the roots of $g_u$.
Signs at $-1,0,1,u,u+1$ give roots $-x,y,\beta$ with
$0<x,y<1<u<\beta<u+1$. Neither rational candidate $\pm1$ is a root.
The norm is $-1$, and the second symmetric function $-1$ gives
$\beta(y-x)=-1+1/\beta<0$. Thus $x>y$ and
$P_n(u)=\lfloor\beta^{3^n}\rfloor$ at every level.
Write
\[
 T(u,w)=(17u+12w+10,24u+17w+15),\qquad
 (a_j,r_j)=T^j(9,14),\quad (b_j,s_j)=T^j(1,4).
\]
The Pell reduction in the supplied round~35 identifies the entire
locus with $a_j$ for $j\ge0$ and $b_j$ for $j\ge1$.
Round~36 excludes the first orbit and all second-orbit indices except
$j\equiv0,4374\pmod{13122}$. We have independently replayed the
complete finite certificate for that assertion. The following
three additional witnesses strengthen its conclusion.

\begin{theorem}\label{thm:pell-new}
If $(P_n(u))_{n\ge0}$ is eventually prime on this Pell locus, then
\[
 u=b_j,\qquad39366\mid j,\qquad
 v_3(u-1)\ge10,\qquad u>10^{59049}.
\]
Every parameter excluded by this condition has a positive lower
density of composite trace indices.
\end{theorem}
\begin{proof}
Put $\eta=17+12\sqrt2$ and write $j=2k$. The identity
\[
 (8b_j+5)+4s_j\sqrt2=(13+16\sqrt2)\eta^j
\]
shows that $\eta^{2m}=1\pmod p$ makes $b_{2k}$ periodic modulo $p$
with period dividing $m$. The following table consists of exact
certificates; its column $c$ is $b_{2a}\bmod p$.
\[
\begin{array}{r|r|r|r|r|r}
p&m&a&c&n&v_3((p-1)(p^2-1)(p^3-1))\\\hline
8279983&6561&2187&7051991&4863973&25\\
27398737&19683&6561&15563564&4929815&31\\
15037811&19683&13122&1777082&549606&9
\end{array}
\]
In every row, $p$ is prime, $\eta^{2m}=1\pmod p$, and
$P_n(c)=0\pmod p$. The root-order return lemma makes each such
zero recurrent, because $n$ exceeds the last column. Thus the entire
parameter-index class $k\equiv a\pmod m$ has infinitely many
composite traces, along an arithmetic progression of trace indices.

For reproducibility, put $M_p=(p-1)(p^2-1)(p^3-1)$ and
$E=3^n\bmod M_p$. Binary polynomial powering modulo $(g_c,p)$ gives
$X^E=A+BX+CX^2$, with coefficient triples in the order of the table
\[
\begin{gathered}
(7459560,4961071,5552814),\qquad (21974350,24684017,2338784),\\
(13021067,3973052,8155287).\end{gathered}
\]
Each satisfies $3A+cB+(c^2+2)C=0\pmod p$. This verifies a large
trace-index witness without constructing the integer trace.
Primality of the three modest witness primes is certified by trial
division up to their square roots.

The round~36 cover leaves $k\equiv0,2187\pmod{6561}$.
The first new row removes the latter class. The two remaining rows
remove $k\equiv6561,13122\pmod{19683}$, leaving $19683\mid k$.
Therefore $39366\mid j$.
The exact identity $v_3(b_j-1)=1+v_3(j)$ for positive even $j$
now gives $v_3(u-1)\ge10$. On the Pell curve, $w>4u/3$, so
$T(u,w)_1>33u$. Hence
\[
 u=b_j>33^j\ge33^{39366}=1089^{19683}>10^{59049}.
\]
\end{proof}

\begin{corollary}[Relative count on the complete Pell locus]
Let $\mathcal P(X)$ count the parameters in the displayed Pell locus
with $u\le X$, and let $\mathcal E(X)$ count those with an eventually
prime trace tower. Then
\[
 \mathcal P(X)=\frac{2\log X}{\log(17+12\sqrt2)}+O(1),\qquad
 \mathcal E(X)\le
 \frac{\log X}{39366\log(17+12\sqrt2)}+O(1).
\]
In particular $\limsup_{X\to\infty}\mathcal E(X)/\mathcal P(X)
\le1/78732$.
\end{corollary}
\begin{proof}
Each affine Pell orbit grows as a positive constant times
$(17+12\sqrt2)^j$. The theorem retains at most every $39366$-th
index of the second orbit and none of the first.
\end{proof}

\begin{remark}[The remaining uniform step]
The strengthened numerical and valuation bounds concern this special
parameter family, not the numerical value of Mills' constant.
The exceptional suborbit is not proved empty. Indeed, for any fixed
modulus $M$, let $L$ be the period of the orbit of $(1,4)$ under $T$
modulo $M$. Such a period exists because the linear part of $T$ has
determinant one. At every positive index divisible by $39366L$,
$b_j\equiv1\pmod M$. The trace pair $(1,-1)$ is an exact fixed point
of the cubing recurrence, so $P_n(b_j)\equiv1\pmod M$ for every $n$.
Thus an arbitrary prescribed finite divisibility sieve still leaves
unboundedly many parameters in the surviving branch.
\end{remark}

% Audited synthesis and additional deductions by local_factors agent.
% Sources: Topic 02 Showa R12,R22,R23,R26,R27,R28; 21 Oct 2026.
% R27 and R28 are already in the supplied archive: do not label them new.
\section{Uniform irreducibility and exact finite divisibility patterns}
\label{sec:factor-patterns}

\paragraph{Provenance and improvement.}
The supplied Showa round~27 already proves uniform irreducibility in
an explicit cyclic cubic unit family. It supersedes the earlier
six-term conditional prefix result as an existence theorem, although
the earlier six polynomials belonged to a different monomial family.
Round~28 already realizes every finite prime/composite pattern under
Schinzel's hypothesis H. We reproduce the decisive proofs and add two
refinements: a general signed-specialization criterion with an exact
three-block limitation, and arbitrary prescribed local multiplicities
which recur unchanged on infinite progressions of trace indices.
These refinements do not remove hypothesis H from any assertion of
simultaneous prime values.

\subsection{The balanced unit and its all-index irreducibility}
Put
\[
\begin{gathered}
A=a^3+2a^2+3a+3,\qquad Q=a^2+a+1,\\
f_a(Z)=Z^3-AZ^2-a^2Z+1,\qquad
d(a)=(a+1)(a^2+3)(a^2+3a+3).\end{gathered}
\]
The discriminant is $d(a)^2$. In the cyclic root ordering set
\[
 s=\frac{ar+1}{1-Qr},\qquad t=\frac{r-1}{Qr+a},
 \qquad\beta_a=-r^5s.
\]
For integer $a\ge8$, take $r$ to be the largest root. The rational-root
test and square discriminant show that the field is cyclic cubic.
The trace seed is
\begin{align*}
 U_0={}&a^{14}+9a^{13}+45a^{12}+161a^{11}+450a^{10}
 +1020a^9+1918a^8\\
 &+3021a^7+3987a^6+4379a^5+3942a^4+2820a^3
 +1521a^2+558a+105,\\
 V_0={}&-3a^6-15a^5-40a^4-75a^3-105a^2-90a-33,
\end{align*}
and
\begin{equation}\label{fac:rec}
 U_{n+1}=U_n^3-3U_nV_n-3,\qquad
 V_{n+1}=V_n^3+3U_nV_n+3.
\end{equation}
These identities follow from $rst=-1$ and Newton's identities:
$U_0=-T_{5,1}$ and $V_0=-T_{4,5}$, where
$T_{ij}=r^is^j+s^it^j+t^ir^j$. Thus
$U_n(a)=\operatorname{Tr}(\beta_a^{3^n})$.

\begin{theorem}[Showa round 27, independently audited]
\label{fac:uniform}
Every $U_n$ is monic and irreducible over $\mathbb Q$, of degree
$14\cdot3^n$, and $U_n(-1)=1$. For every integer $a\ge8$ the number
$\beta_a$ is a cubic Pisot unit of norm $-1$ whose conjugates are
$\beta_a,-X,Y$, with
\[
 0<Y<X<1,\quad X^{40}\beta_a^{17}<1,
 \quad X^{23}<Y^{17},\qquad
 U_n(a)=\lfloor\beta_a^{3^n}\rfloor.
\]
\end{theorem}
\begin{proof}
First every odd trace is positive at every real parameter. For real
$a\ne-1$, the three roots of $f_a$ are distinct and real. Their product
is $-1$ and their pairwise sum is $-a^2\le0$, so they have the form
$p>q>0,-w<0$ with $pqw=1$. If $a>-1$, then $f_a(1)<0$ gives
$q<1<p$. The cyclic fractional-linear map sends $p$ to $-w$:
its value cannot vanish or have a pole on this connected parameter
interval, since its determinant is $(a+1)^2\ne0$, and it is negative
at $a=0$. The conjugates of $\beta_a$ are therefore
\[
 P=p^5w,\qquad R=w^5q,\qquad -B=-q^5p,
 \qquad B/P=q^6/p^3<1.
\]
Their odd-power sum is positive. If $a<-1$, then $f_a(1)>0$ and
$A(a)<A(-1)=1$ show that both positive roots lie below one: if both
exceeded one, $A=p+q-1/(pq)>1$. In either cyclic orientation the
negative conjugate has magnitude $B<1$, while $PRB=1$ forces a positive
conjugate above one. The odd trace is again positive. At $a=-1$ the
specialized conjugate multiset is $(1,1,-1)$, and every odd trace is one.

Direct reduction gives
\begin{equation}\label{fac:res}
 \overline U_0=a^5(a-1)^3(a^2+1)^3\quad\text{over }\mathbb F_3.
\end{equation}
For $n\ge1$ the three Hensel blocks of $U_n$ are irreducible over
$\mathbb Q_3$ and have degrees $5\cdot3^n,3\cdot3^n,6\cdot3^n$.
Here is the valuation argument, including repeated residual factors.
At a root $\xi$ of $U_n$, monicity gives integrality and
\[
 v_3(U_j(\xi))=3^{j-n}\quad(0\le j<n).
\]
The last step follows from $U_{n-1}^3=3(1+U_{n-1}V_{n-1})$;
descending steps compare valuations less than one with terms divisible
by three. If a residual irreducible factor $h$ has multiplicity $m$
and degree $e$, write $U_0=H^m B+3C$ with $B(\xi)$ a unit.
Then $v_3(H(\xi))=1/(m3^n)$. The ramification index is divisible by
$m3^n$ and the residue degree by $e$. Their product equals the degree
of the Hensel block, proving its irreducibility.

A monic rational factor $F$ of $U_n$ is integral and has no real root,
so it has even degree and is positive on the real line. The only proper
even subset sums of $5,3,6$ are $6$ and $8$. The corresponding residual
factors are
\[
 (a^2+1)^{3^{n+1}},\qquad
 a^{5\cdot3^n}(a-1)^{3^{n+1}}.
\]
Both have value $-1$ modulo three at $a=-1$. On the other hand $F$ and
its monic complement have positive integer values there multiplying
to $U_n(-1)=1$, so $F(-1)=1$. This contradiction proves irreducibility
for every $n\ge1$.

For the seed, the exact factorizations modulo two have degrees $6,8$,
and those modulo thirteen have degrees $5,9$. For reproducibility,
the factors are specified by ascending coefficient arrays:
\[
\begin{array}{c|l}
2&(1,0,1,1,0,1,1),\quad(1,0,0,1,0,1,1,0,1)\\
13&(12,7,8,8,4,1),\quad(12,7,2,11,12,9,6,4,5,1).
\end{array}
\]
Their products and irreducibility are certified by the Frobenius
criterion $X^{p^d}=X\bmod F$ and
$\gcd(F,X^{p^{d/q}}-X)=1$ for primes $q\mid d$. The disjoint proper
factor-degree sets prove rational irreducibility of $U_0$.

Finally, for $a\ge8$ and $w=-s$,
\[
 a^3<r<2a^3,\quad a/Q<w<1/a,\quad 0<t<1/a^2,
 \quad rwt=1,
\]
so
\[
 (4/5)a^{14}<\beta_a<32a^{14},\quad
 X=t^5r<2a^{-7},\quad Y=w^5t<a^{-7}.
\]
Since $V_0(a)<-1$ and $\beta_aXY=1$,
$\beta_a(Y-X)=V_0(a)+1/\beta_a<0$, giving $X>Y$.
Also $X^{40}\beta_a^{17}<2^{125}a^{-42}\le1/2$, and
$X^{23}/Y^{17}=X^{40}\beta_a^{17}$. Every odd-power trace is its
positive dominant power minus a number strictly between zero and one,
which proves the floor identities.
\end{proof}

\subsection{A signed-specialization criterion and its exact limitation}
\begin{proposition}\label{fac:signcriterion}
Suppose integer polynomials $U_0,V_0$ satisfy \eqref{fac:rec}, $U_0$ is
monic, $\deg V_0<2\deg U_0$, and for some integer $a_*$ one has
$(U_0(a_*),V_0(a_*))=(1,-1)$. Suppose each $U_n$ is strictly positive
on $\mathbb R$. Write
$\overline U_0=\prod_{i=1}^r h_i^{m_i}$ into distinct monic irreducible
factors over $\mathbb F_3$, with $\deg h_i=d_i$. Give block $i$ the label
\[
 \lambda_i=(m_id_i\bmod2,\ h_i(a_*)^{m_i})
 \in(\mathbb Z/2\mathbb Z)\times\mathbb F_3^\times.
\]
If no nonempty proper subcollection of these labels has product
$(0,1)$, then $U_n$ is irreducible over $\mathbb Q$ for every $n\ge1$.
Moreover, this label criterion is possible only for $r\le3$. With
$r=3$, it holds exactly when the labels are the three different
nonidentity elements of this group.
\end{proposition}
\begin{proof}
The root-valuation proof above gives one irreducible local block of
degree $m_id_i3^n$ and reduction $h_i^{m_i3^n}$ for each $i$.
Since $3^n$ is odd, its parity and value at $a_*$ retain precisely the
label $\lambda_i$. A proper rational factor must have even degree and
positive integer value one at $a_*$, contradicting the label condition.
The product of all labels is $(0,1)$, because $U_0$ has even degree
and value one at $a_*$. The group has three nonidentity elements,
each of order two. A minimal product-one sequence cannot contain the
identity (unless it has length one), or a repeated element (unless
it has length two). Thus it has length at most three; the three-element
case must consist of all three nonidentity elements. Conversely those
three labels have product one and no proper product-one subcollection.
\end{proof}
This is a limitation of this particular factor obstruction, not a
bound on the number of local factors of a globally irreducible polynomial.
For the balanced seed the labels are $(1,-1),(1,+1),(0,-1)$, so the
criterion is extremal. The four-block degree-$31$ seed of the earlier
six-prefix argument cannot be handled by this exact criterion alone.

\subsection{Simple forcing roots and exact local multiplicities}
We next state a strengthening of round~28. All prime-value assertions
below explicitly assume the usual integer-prime hypothesis H.

\begin{lemma}[Forcing roots, from round 28]\label{fac:forcing}
For every finite window $I=\{0,\ldots,L\}$ and $j\in I$, a positive-
Dirichlet-density set of primes $p\equiv2\pmod3$ has a residue $b$ such
that
\[
 U_j(b)=0,\quad U'_j(b)\ne0,\quad d(b)\ne0,
 \qquad U_i(b)\ne0\quad(i\in I\setminus\{j\}).
\]
Any prescribed finite set of primes may be excluded.
\end{lemma}
\begin{proof}
Let $K_j=\mathbb Q(\xi)$ for a root of $U_j$. It has an odd-degree
completion above three. For $j\ge1$ this follows from the local blocks
of degrees $5\cdot3^j$ and $3\cdot3^j$. For $j=0$, Hensel lifts the
residual factor $a^5$ to a degree-five block; at least one of its
irreducible constituents has odd degree. Thus
$\mathbb Q(\sqrt{-3})$ is not a subfield of $K_j$, since its completion
at three has degree two and would embed in every completion of $K_j$.

In a Galois extension containing the splitting field of $U_j$ and
$\mathbb Q(\sqrt{-3})$, choose an automorphism fixing $\xi$ and acting
nontrivially on the quadratic field. It exists by the preceding
noncontainment. Its conjugacy class has a fixed root and the same
nontrivial quadratic character. Chebotarev gives a positive-density
set of primes with a root of $U_j$ and $p\equiv2\pmod3$.
Exclude the finitely many divisors of $\operatorname{disc}(U_j)$,
$\operatorname{Res}(U_j,d)$, and
$\operatorname{Res}(U_j,U_i)$ for $i\ne j$. All these integers are
nonzero by irreducibility and unequal degrees. The asserted simple
root and all other nonvanishing conditions follow.
\end{proof}

\begin{theorem}[Prescribed exact multiplicities with periodic recurrence]
\label{fac:valuationpatterns}
Fix a finite window $I=P\sqcup C$, a positive integer $M$, and for each
$j\in C$ a nonempty finite list of positive integers
$e_{j,1},\ldots,e_{j,r_j}$. There are pairwise distinct primes
$p_{j,k}\ge5$ with $p_{j,k}\equiv2\pmod3$, all coprime to $M$, and an arithmetic progression
$\mathcal A$ of integer parameters with the following unconditional
properties for every sufficiently large $a\in\mathcal A$:
\begin{enumerate}
\item $a\equiv-1\pmod M$, and every odd-power trace of $\beta_a$ is
$1$ modulo $M$;
\item for $i\in I$,
\[
 v_{p_{j,k}}(U_i(a))=
 \begin{cases}e_{j,k},&i=j,\\0,&i\ne j;\end{cases}
\]
\item with
\[
 T_{j,k}=\operatorname{ord}_{\,p_{j,k}^{e_{j,k}}
                      (p_{j,k}^3-1)}(3),
\]
one has
\[
 v_{p_{j,k}}(U_{j+hT_{j,k}}(a))=e_{j,k}\qquad(h\ge0).
\]
\end{enumerate}
Assuming hypothesis H, infinitely many $a\in\mathcal A$ additionally
make every $U_i(a)$, $i\in P$, prime. Every prescribed composite position
$j$ then has at least $r_j$ different prime divisors, with exactly the
specified multiplicities for the chosen divisors. If $P$ is empty,
no conjecture is required.
\end{theorem}
\begin{proof}
For each pair $(j,k)$ choose a fresh prime $p\ge5$ and simple forcing root
$(p,b)$ from Lemma~\ref{fac:forcing}. Hensel's lemma lifts $b$ uniquely
to a root modulo $p^{e+1}$, where $e=e_{j,k}$. Among the $p$ lifts of
the root modulo $p^e$, choose any of the other $p-1$ classes modulo
$p^{e+1}$. Call it $c$. Taylor expansion and $U'_j(b)\ne0$ give
$v_p(U_j(c))=e$; every other designated $U_i(c)$ is a unit modulo $p$.
CRT now imposes $a\equiv-1\pmod M$ and
$a\equiv c_{j,k}\pmod{p_{j,k}^{e_{j,k}+1}}$ simultaneously. The resulting
progression has modulus
\[
 D=M\prod_{j\in C}\prod_{k=1}^{r_j}p_{j,k}^{e_{j,k}+1}.
\]
This proves the exact multiplicities within the finite window.

We verify recurrence at the same precision. Fix a chosen $p,e$ and
consider
\[
 \mathcal B=(\mathbb Z/p^{e+1}\mathbb Z)[Z]/(f_a(Z)).
\]
Because $p\nmid d(a)$, its reduction is etale. The square discriminant
forces splitting type $1+1+1$ or $3$, so every residue-field unit is
killed by $p^3-1$. The kernel of reduction on $\mathcal B^\times$ is
$1+p\mathcal B$ and is killed by $p^e$, by induction with the binomial
theorem. Hence every unit of $\mathcal B$ is killed by
$E=p^e(p^3-1)$. The root transformation defining $s$ has invertible
denominator: a simultaneous zero of $1-QZ$ and $aZ+1$ would force
$(a+1)^2=0\bmod p$, excluded by $d(a)\ne0$. Thus $\beta=-r^5s$ is a unit
in this algebra, and its traces agree with the integer polynomials.
Since $p\equiv2\pmod3$, $\gcd(E,3)=1$. For $T=\operatorname{ord}_E(3)$,
\[
 U_{n+T}(a)\equiv U_n(a)\pmod{p^{e+1}}\qquad(n\ge0).
\]
The nonzero residue of $U_j(a)/p^e$ modulo $p$ is therefore retained
at every $j+hT$, proving the exact recurring valuations.

Write the CRT progression as $a=A_0+Dz$. For each $i\in P$, the
polynomial $U_i(A_0+Dz)$ is irreducible, primitive, and has positive
leading coefficient. Their product has no fixed prime divisor:
primes dividing $M$ give values one; each forcing prime gives nonzero
values at every prime-designated index; at a prime not dividing $D$
choose $z$ so that $A_0+Dz\equiv-1$. Hypothesis H consequently supplies
infinitely many simultaneous prime values. All other claims are
unconditional congruences. Finally, every odd trace is an integer
polynomial in $a$ and has value one at $a=-1$, proving the modulo-$M$
assertion. Taking $a$ large ensures the Pisot bounds and makes each
forced value larger than its chosen prime divisors.
\end{proof}

The theorem prescribes valuations of selected divisors, not a complete
prime factorization: unselected prime factors may also occur. It
strengthens the finite-pattern statement in the archive while retaining
its essential finite-window quantifier and its explicit conjectural
prime-value assumption.

\subsection{Audit boundaries and remaining arithmetic question}
The single-block norm-minus-one family of round~22 also has all-index
local irreducibility, but its index-zero trace is a fourth power; its
prime-prefix statement deliberately begins at the next index. The
cyclic unit family above is stronger for conditional prime-prefix
existence and starts at index zero. Round~23 is compatible with the
present recurrence theorem: its isolated divisors are chosen at primes
congruent to one modulo three, while the recurring divisors above lie
at primes congruent to two modulo three. In round~23 the prime and unit
may depend on the isolated index; the result does not say that every
trace has such a divisor in one fixed tower.

Round~26 proves a separate polynomial-family obstruction: fixed nonzero
norm and square discriminant force the polynomial coefficient functions
to be algebraically dependent and, over $\mathbb C$, to factor through
one common polynomial. Its logarithmic two-form argument and the
normalization-to-an-affine-line argument are valid; they do not prevent
the one-parameter construction used here. We do not use a conjectural
cyclic lift of the round~22 family, the multivariate Hilbert construction
of round~24, or an unverified primality certificate from round~25 as
inputs to the proofs above.

Uniform irreducibility is an unconditional result already in the
archive. Arbitrarily long prime prefixes remain conditional on H;
exact prescribed local divisibility is unconditional. None implies
one fixed algebraic number whose entire infinite trace tower is prime.
Nor does any identify the least-prime path defining Mills' constant.
The full cubic Pisot compositeness problem and unconditional Mills
transcendence remain open within these arguments.

% Transfer audit and a quantitative ray-unit barrier.
% Written 21 October 2026. This is a section fragment, not a standalone file.
% New relative to the audited older [031] and supplied non-Topic02 notes;
% no literature-wide priority assertion is intended.

\section{Finite local conditions and logarithmic unit lattices}
\label{transfer:sec:rayunits}

The new non-Mills material does not supply an additional proved route to
Mills transcendence. A useful transfer is nevertheless possible: the
unit-lattice viewpoint gives a quantitative explanation of why finitely
many local conditions and Saito's strict conjugate bounds can coexist.
It identifies a genuine limitation of that combination of tests without
assuming that any trace is prime.

Let $K$ be a totally real cubic field, with real embeddings
$\sigma_1,\sigma_2,\sigma_3$. Fix a unit $w\in\mathcal O_K^\times$ with
\[
 \sigma_1(w)>0,\qquad \sigma_2(w)<0,\qquad \sigma_3(w)>0.
\]
Its norm is $-1$. Let $\mathfrak m$ be a nonzero integral ideal, and set
\[
 E^+(\mathfrak m)=
 \{u\in\mathcal O_K^\times:
       \sigma_i(u)>0\ (1\le i\le3),\ u\equiv1\pmod{\mathfrak m}\}.
\]
Its two-coordinate logarithmic lattice is
\[
 \Lambda_{\mathfrak m}=
 \{(\log\sigma_1(u),\log\sigma_2(u)):u\in E^+(\mathfrak m)\}
 \subset\mathbb R^2.
\]
Write $\Delta_{\mathfrak m}>0$ for its covolume in these coordinates.
This convention avoids an extraneous factor of $\sqrt3$ from embedding
the lattice into the plane $x+y+z=0$ in $\mathbb R^3$.

\begin{theorem}[Quantitative persistence in every admissible unit coset]
\label{transfer:thm:ray}
Fix $0<a<1/2$. For $H>1$, let $B_{w,\mathfrak m,a}(H)$ count the units
$\beta\in wE^+(\mathfrak m)$ such that, on writing
\[
 B=\sigma_1(\beta),\qquad r=\sigma_2(\beta),\qquad
 s=\sigma_3(\beta),
\]
one has
\[
 1<B\le H,\qquad 0<s<-r<1,\qquad -r<B^{-a}.
\]
Then
\begin{equation}\label{transfer:eq:count}
 B_{w,\mathfrak m,a}(H)
 =\frac{1-2a}{4\Delta_{\mathfrak m}}(\log H)^2
     +O_{K,w,\mathfrak m,a}(\log H+1).
\end{equation}
Every counted $B$ is a cubic Pisot number of norm $-1$. In addition,
\[
 -r<s^{a/(1-a)},\qquad \beta\equiv w\pmod{\mathfrak m}.
\]
In particular, for $a=17/40$, the two strict bounds
\[
 -r<\min(B^{-17/40},s^{17/23})
\]
hold for quadratically many units in logarithmic height in every such
coset.
\end{theorem}

\begin{proof}
The subgroup $E^+(\mathfrak m)$ has finite index in the unit group:
positivity imposes only a finite sign condition, and reduction modulo
$\mathfrak m$ maps units to the finite group
$(\mathcal O_K/\mathfrak m)^\times$. Dirichlet's unit theorem therefore
makes $\Lambda_{\mathfrak m}$ a full lattice of rank two. Positivity
removes its possible torsion, so its logarithm map is injective.

The logarithms of the coset $wE^+(\mathfrak m)$ form the translate
\[
 (\log\sigma_1(w),\log|\sigma_2(w)|)+\Lambda_{\mathfrak m}.
\]
For a point of this translate put $x=\log B$ and $y=\log(-r)$.
Because $|\operatorname N\beta|=1$, one has $\log s=-x-y$.
The asserted inequalities are exactly
\begin{equation}\label{transfer:eq:triangle}
 0<x\le\log H,\qquad -x/2<y<-ax.
\end{equation}
Indeed $y>-x/2$ is $-r>s$, and $y<-ax<0$ gives $-r<1$;
then $-x-y<-x/2<0$, so $s<1$ as well.
The triangle in \eqref{transfer:eq:triangle} has area
\[
 \int_0^{\log H}(1/2-a)x\,dx
       =\frac{1-2a}{4}(\log H)^2
\]
and boundary length $O_a(\log H)$. Counting translates of a fixed
fundamental parallelogram gives its area divided by
$\Delta_{\mathfrak m}$, with error $O(\log H+1)$: only cells meeting a
fixed-width neighborhood of the boundary can contribute to the error.
Strict versus non-strict boundaries changes the count by at most the
same order. This proves \eqref{transfer:eq:count}.

The dominant embedding exceeds one and the other two have modulus
less than one, so $B$ is Pisot. It cannot be rational, since a rational
unit is $\pm1$; its degree is therefore three. Finally
$B(-r)s=1$, and the inequality $-r<B^{-a}$ is equivalent to
$(-r)^{1-a}<s^a$, proving the additional bound. The congruence is the
one defining the coset.
\end{proof}

\begin{corollary}[Persistence of every fixed finite trace sieve]
\label{transfer:cor:sieve}
Take $\mathfrak m=M\mathcal O_K$ with integer $M\ge2$. Every unit
counted in Theorem~\ref{transfer:thm:ray} satisfies
\[
 \operatorname{Tr}_{K/\mathbb Q}(\beta^j)
 \equiv\operatorname{Tr}_{K/\mathbb Q}(w^j)\pmod M
 \qquad(j\ge1).
\]
Thus, if a seed $w$ has
$\gcd(\operatorname{Tr}(w^{3^n}),M)=1$ at every level $n\ge0$,
then all the counted balanced Pisot units have the same avoidance
property at every level.
\end{corollary}
\begin{proof}
From $\beta-w\in M\mathcal O_K$ one obtains
$\beta^j-w^j\in M\mathcal O_K$ for every positive integer $j$.
The trace of an algebraic integer is an integer, so tracing gives the
stated congruence. The last assertion follows separately at every
level; no primality assertion is inferred from coprimality with $M$.
\end{proof}

The seed hypothesis in this corollary is essential. The theorem does
not assert that an arbitrary finite sieve has an avoiding seed in a
specified field. Instead it proves that, once a finite residue class
is admissible, the archimedean restrictions alone cannot eliminate it:
there are asymptotically $c(\log H)^2$ members satisfying the strict
conjugate bounds. The regulator enters only through the positive
covolume $\Delta_{\mathfrak m}$ and changes the leading constant,
not the existence of these candidates. This is stronger than an
unquantified ray construction but does not close the remaining
norm-minus-one cyclic Pisot case.


\begin{example}[An explicit nonempty congruence coset]
\label{transfer:ex:seed}
Let $r$ be the largest root of $X^3-25X^2-4X+1$ and order its other
roots as $s<0<t$. The rational root test excludes $\pm1$, the
discriminant is $273^2$, and root brackets give
\[
 25<r<26,\qquad -0.294<s<-0.292,\qquad 0.135<t<0.137.
\]
Thus $K=\mathbb Q(r)$ is cyclic cubic and $r$ has signs $(+,-,+)$
and norm $-1$. Direct root enumeration of the polynomial over
$\mathbb F_2,\mathbb F_5,\mathbb F_{17}$ shows it irreducible in all
three cases; these primes are inert. The integral element
$v=(r^2+2)/3$ satisfies
\[
 v^3-213v^2+290v-99=0.
\]
This polynomial reduces to $X(X-1)(X+1)$ modulo three, so $3$ splits
completely. The residues of $r$ at its three primes are
$(-1,1,1)$ in some order, as follows from its norm $-1$ and trace
$25\equiv1\pmod3$.

Put $w=r^{66619}$, where $66619=7\cdot31\cdot307$.
At $2$, the order of $r$ divides $7$; at $5$ and $17$, the finite-field
norm identities give $r^{31}=-1$ and $r^{307}=-1$, respectively.
The relevant quotients of $66619$ are odd. Consequently, for every
odd integer $j\ge1$,
\[
 \operatorname{Tr}(w^j)\equiv
 \begin{cases}
 1&\pmod2,\\
 1&\pmod3,\\
 2&\pmod5,\\
 14&\pmod{17}.
 \end{cases}
\]
Applying Theorem~\ref{transfer:thm:ray} to
$wE^+(510\mathcal O_K)$ gives asymptotically
$3(\log H)^2/(80\Delta_{510\mathcal O_K})$ strictly Saito-balanced
Pisot units whose entire cubic-power trace towers avoid
$2,3,5,17$. This asserts finite-prime avoidance, not primality.
\end{example}

% Section fragment for the integrated revision; written 21 October 2026.
% Uses thm:windows, lem:contraction, prop:fourteen, Keio009, and S14.
\section{A complete cyclic cubic norm class}
\label{rescyc:sec:cyclic}
Submitted round S14 already proves that an eventually prime,
growing trace tower in a cyclic cubic field must have norm congruent
to zero or minus one modulo three \cite[Theorem~1]{S14}. Thus the
norm class excluded below is not newly discovered here. Combining
that restriction with the trace-limit classification and cyclic
Frobenius-return argument of \cite{Keio009}, we record a
positive-density formulation for the actual Pisot floors, including
both possible real floor branches. This is a synthesis of existing
ingredients, not a new return mechanism or a claim of historical
priority.

\begin{theorem}[Cyclic cubic norms congruent to one]
\label{rescyc:thm:norm-one-mod-three}
Let $K/\Q$ be a cyclic cubic extension, and let $\beta\in K$ be a
cubic Pisot number. If
\[
 N=\Nm_{K/\Q}(\beta)\equiv1\pmod3,
\]
then the set
\[
 \mathcal C_\beta=
 \{n\ge0:\lfloor\beta^{3^n}\rfloor\text{ is composite}\}
\]
has positive lower density. In particular, this conclusion holds
for every cyclic cubic Pisot unit of norm $+1$.
\end{theorem}

\begin{proof}
A cyclic cubic field is totally real, since a nontrivial complex
conjugation would have order two in a group of order three. Write
the conjugates of $\beta$ as $\beta,\gamma_2,\gamma_3$, and put
\[
 U_n=\Tr_{K/\Q}(\beta^{3^n}),\qquad
 f(X)=X^3-uX^2+vX-N.
\]
The two small conjugates have different absolute values: equality
would give $\gamma_2=-\gamma_3$ and hence $\beta=u\in\Z$,
contrary to degree three. Thus their odd-power sum has a fixed
nonzero sign. For all sufficiently large $n$,
\[
 \lfloor\beta^{3^n}\rfloor=U_n-s,
 \qquad s\in\{0,1\},
\]
where $s=1$ exactly when the small conjugate of larger absolute
value is positive. Also $U_n\to+\infty$.
Let $\lambda=\lim_{n\to\infty}U_n$ in $\Z_3$.

If $s=1$, Proposition~\ref{prop:fourteen} and its explicit table
show that either equality $\lambda=0$ or $\lambda=2$ would force
$3\mid N$. Since $N\equiv1\pmod3$, neither equality is possible.
Theorem~\ref{thm:windows}, with the single shift $h=-1$, supplies
an eventual arithmetic progression of composite floors.

Suppose now that $s=0$. If $\lambda\notin\{1,-1\}$,
Theorem~\ref{thm:windows} with $h=0$ gives the same conclusion.
We handle the two remaining limits separately.

If $\lambda=1$, the residue classification gives
\[
 f(X)\equiv(X-1)(X^2+1)\pmod3,
\]
because $N\equiv1\pmod3$. Its root $1$ is simple, so Hensel's
lemma gives an embedding $K=\Q(\beta)\hookrightarrow\Q_3$.
In a Galois extension all primes over three have the same local
degree; the degree-one completion therefore implies that three
splits completely in $K$. The three conjugates of $\beta$ now lie
in $\Q_3$ and are units, because $3\nmid N$. Each has
Teichm\"uller part $1$ or $-1$. Their sum is $\lambda=1$, so
these three parts must be $1,1,-1$ in some order. Their product
gives $N\equiv-1\pmod3$, a contradiction. Equivalently, the
displayed separable reduction has factorization type $1+2$,
which cannot occur at an unramified prime in a cyclic cubic field.

It remains to treat $\lambda=-1$. Lemma~\ref{lem:contraction}
gives $U_n\equiv-1\pmod3$ for every $n$. If only finitely many
$U_n$ are prime, then every sufficiently late floor is composite,
and the density is one. Otherwise choose a sufficiently late
prime $p=U_r$ with
\[
 p\nmid3N\operatorname{disc}(f).
\]
This is possible since $U_n\to+\infty$. Frobenius in the cyclic
cubic field has cycle type $1+1+1$ or $3$. Consequently all roots
of $f$ modulo $p$ lie in $\F_{p^3}$ and are nonzero. Since
$p\equiv-1\pmod3$,
\[
 \gcd(3,p^3-1)=1.
\]
Choose $D\ge1$ with $3^D\equiv1\pmod{p^3-1}$. For every
residual root $a$ and every $j\ge0$ this gives
$a^{3^{r+jD}}=a^{3^r}$. Summing over the three roots yields
\[
 U_{r+jD}\equiv U_r=0\pmod p\qquad(j\ge0).
\]
For sufficiently large $j$ the floor equals $U_{r+jD}>p$, so
this is an arithmetic progression of composite floors. Every
progression obtained above has positive lower density, completing
the proof.
\end{proof}

\begin{remark}[What remains in the cyclic Mills branch]
For a hypothetical cyclic cubic shadow of the least Mills
constant, the dominant small conjugate is negative. An eventually
prime trace tail therefore forces $\lambda=1$: other limits are
excluded by Theorem~\ref{thm:windows} or by the last part of the
proof above. If its norm is prime to three, complete splitting at
three then forces $N\equiv-1\pmod3$. In particular, a unit
shadow must have norm $-1$. The norm-$-1$ cyclic units constructed
elsewhere in this paper show why Theorem~\ref{rescyc:thm:norm-one-mod-three}
does not settle the entire cyclic case or unconditional Mills
transcendence.
\end{remark}


\section{Excluding rational powers of Mills' constant from number fields}
\label{res:sec:fields}

Restricted cubic results already give complete exclusion statements for
other, substantially larger classes of algebraic expressions. Throughout
this section $A$ is the least Mills constant defined in the introduction.
The cubic three-adic embedding obstruction below was already proved in
the supplied manuscript \cite{Keio031}. We include a short proof from
Theorem~\ref{thm:windows} and Proposition~\ref{prop:fourteen} to make the
subsequent field deductions self-contained. The ambient-field and
cyclotomic consequences are deductions from this obstruction; no
historical priority is asserted for them.

\begin{lemma}[Prime traces force a three-adic embedding]
\label{res:lem:prime-embedding}
Let $\beta$ be a cubic Pisot number, with minimal polynomial
$f(X)=X^3-uX^2+vX-N$, and put
$U_n=\Tr_{\Q(\beta)/\Q}(\beta^{3^n})$.
If $U_n$ is prime for every sufficiently large $n$, then its three-adic
limit satisfies $\lambda\in\{1,-1\}$, and $\Q(\beta)$ admits an
embedding into $\Q_3$. No restriction on the norm $N$ is required.
\end{lemma}
\begin{proof}
The Pisot property gives $U_n=\beta^{3^n}+o(1)$, so $U_n\to\infty$.
If $\lambda\notin\{1,-1\}$, apply Theorem~\ref{thm:windows} with the
single shift $0$. It gives infinitely many composite $U_n$, contrary
to the hypothesis. Write $\lambda=\varepsilon$, where
$\varepsilon\in\{1,-1\}$. The six corresponding coefficient classes
in Proposition~\ref{prop:fourteen} satisfy
\[
 u\equiv\varepsilon,\qquad N\equiv\varepsilon v\pmod3,
 \qquad
 f(X)\equiv(X-\varepsilon)(X^2+v)\pmod3.
\]
If $v\not\equiv2\pmod3$, the residual root $\varepsilon$ is simple.
If $v\equiv2\pmod3$, the residual polynomial is
$(X-\varepsilon)^2(X+\varepsilon)$, whose root $-\varepsilon$ is simple.
Hensel's lemma therefore gives a root of $f$ in $\Z_3$. Since $f$ is
irreducible over $\Q$, evaluation at this root embeds
$\Q[X]/(f)=\Q(\beta)$ into $\Q_3$.
\end{proof}

\begin{lemma}[Pisot powers retain their field]
\label{res:lem:power-field}
If $\beta$ is a Pisot number and $r$ is a positive integer, then
$\Q(\beta^r)=\Q(\beta)$.
\end{lemma}
\begin{proof}
An embedding of $\Q(\beta)$ into $\mathbb C$ that fixes
$\Q(\beta^r)$ sends $\beta$ to a conjugate $\gamma$ with
$\gamma^r=\beta^r$. Hence $|\gamma|=\beta>1$. The defining Pisot
condition leaves only $\gamma=\beta$. Separability in characteristic
zero says that the number of such embeddings is
$[\Q(\beta):\Q(\beta^r)]$, which is therefore one.
\end{proof}

\begin{theorem}[An obstruction in every ambient number field]
\label{res:thm:ambient}
Let $K\subseteq\mathbb C$ be a number field. If $A^q\in K$ for some
$q\in\Q\setminus\{0\}$, then $K$ contains a totally real cubic
subfield $F$ admitting an embedding into $\Q_3$. Consequently, if
$K$ has no such subfield, then
\[
 \boxed{A^q\notin K\qquad(q\in\Q\setminus\{0\}).}
\]
Here rational powers of $A>1$ denote their positive real values.
\end{theorem}
\begin{proof}
Replacing $q$ by $-q$ if necessary, write $q=r/s$ with positive
integers $r,s$. The inclusion $A^{r/s}\in K$ makes $A$ algebraic.
Saito's reduction \cite[Theorem~1.7]{SaitoVariants} supplies
$\beta=A^{3^m}$, a totally real cubic Pisot number satisfying
\eqref{eq:saitobounds}. Its two small conjugates have negative sum
at every sufficiently large odd power: the negative one has strictly
larger absolute value. Thus
\[
 \lfloor\beta^{3^n}\rfloor
 =\Tr_{\Q(\beta)/\Q}(\beta^{3^n})
\]
for every sufficiently large $n$. These integers are prime by the
definition of $A$. Lemma~\ref{res:lem:prime-embedding} therefore
gives a $\Q_3$-embedding of $F=\Q(\beta)$. Finally,
\[
 \beta^r=(A^{r/s})^{s3^m}\in K,
 \qquad \Q(\beta^r)=F
\]
by Lemma~\ref{res:lem:power-field}, proving $F\subseteq K$.
\end{proof}

\begin{corollary}[Cubic fields and radical expressions]
\label{res:cor:cubic}
If a cubic number field $K$ has no embedding into $\Q_3$, then no
nonzero rational power of $A$ belongs to $K$. In particular this holds
when $3$ is inert or totally ramified in $K$. Equivalently, no positive
element of such a field, raised to a nonzero rational power, can
equal $A$.
\end{corollary}
\begin{proof}
Any cubic subfield of the cubic field $K$ is $K$ itself, so
Theorem~\ref{res:thm:ambient} applies. Inertness and total ramification
both exclude a prime of $K$ over $3$ with ramification index and
residue degree equal to one, which an embedding into $\Q_3$ would
require.
\end{proof}

\begin{theorem}[A Galois prime-count obstruction]
\label{res:thm:galois}
Let $K/\Q$ be a finite Galois extension and let $g_3(K)$ be the
number of prime ideals of $K$ above $3$. If $3\nmid g_3(K)$, then
$A^q\notin K$ for every $q\in\Q\setminus\{0\}$.
\end{theorem}
\begin{proof}
Suppose otherwise and let $F\subseteq K$ be the cubic subfield from
Theorem~\ref{res:thm:ambient}. Put
$G=\operatorname{Gal}(K/\Q)$ and $H=\operatorname{Gal}(K/F)$, so
$[G:H]=3$. The $\Q_3$-embedding supplies a prime $\mathfrak p$ of
$F$ with completion $F_{\mathfrak p}=\Q_3$. Choose a prime
$\mathfrak P$ of $K$ above $\mathfrak p$, and let $D\subseteq G$ be
its decomposition group. The image of $F$ in $K_{\mathfrak P}$ lies
in $\Q_3$, so every member of $D$ fixes $F$ and $D\subseteq H$.
The usual transitive action on the primes above $3$ gives
\[
 g_3(K)=[G:D]=[G:H][H:D]=3[H:D],
\]
contrary to the hypothesis.
\end{proof}

\begin{remark}[The exact abelian scope of the local obstruction]
\label{res:rem:abelian}
If $K/\Q$ is finite abelian, it contains a cyclic cubic subfield in
which $3$ splits completely if and only if $3\mid g_3(K)$.
Indeed, its decomposition group $D$ is normal. An index-three
subgroup $H$ containing $D$ exists precisely when the finite abelian
group $G/D$ has order divisible by three. Its fixed field is cyclic
cubic, is totally real because its degree is odd, and has trivial
decomposition group at $3$. This equivalence concerns the existence
of a subfield with the required local property; it does not assert
that a Mills candidate exists in that subfield.
\end{remark}

\begin{corollary}[Cyclotomic exclusions]
\label{res:cor:cyclotomic}
Write $N=3^b m$, where $b\ge0$ and $(m,3)=1$. If
\[
 3\nmid\frac{\varphi(m)}{\ord_m(3)},
\]
then $A^q\notin\Q(\zeta_N)$ for every nonzero rational $q$.
Here $\zeta_N$ is a primitive $N$th root of unity, and
$\varphi(1)=\ord_1(3)=1$ by convention. In particular, for every
nonzero rational $q$,
\[
 A^q\notin\bigcup_{a,b\ge0}\Q(\zeta_{7^a3^b}),
 \qquad
 A^q\notin\bigcup_{a,b\ge0}\Q(\zeta_{13^a3^b}).
\]
\end{corollary}
\begin{proof}
In $\Q(\zeta_m)$ the prime $3$ is unramified and its Frobenius has
order $\ord_m(3)$, so it has
$\varphi(m)/\ord_m(3)$ primes above it. The $3$-power cyclotomic
factor is totally ramified at $3$. Locally, adjoining this totally
ramified factor to the unramified completion multiplies the local
degree by $\varphi(3^b)$, just as the global degree is multiplied
by $\varphi(3^b)$. Thus it does not change the number of primes:
\[
 g_3\bigl(\Q(\zeta_{3^b m})\bigr)
 =\frac{\varphi(m)}{\ord_m(3)}.
\]
Theorem~\ref{res:thm:galois} proves the first assertion.

For $a\ge1$, elementary lifting of multiplicative orders gives
\[
 \ord_{7^a}(3)=6\cdot7^{a-1},\qquad
 \ord_{13^a}(3)=3\cdot13^{a-1}.
\]
Indeed, the orders modulo $7$ and $13$ are $6$ and $3$, respectively,
and $v_7(3^6-1)=v_{13}(3^3-1)=1$. The corresponding values of
$g_3$ are $1$ and $4$, neither divisible by three. The case $a=0$
has $g_3=1$. Every element of either displayed union lies in one
finite stage, proving the two exclusions.
\end{proof}

\begin{corollary}[A quotient obstruction]
\label{res:cor:quotient}
Let $K/\Q$ be finite Galois. If
$\operatorname{Gal}(K/\Q)$ has no quotient isomorphic to $C_3$ or
$S_3$, then $A^q\notin K$ for every nonzero rational $q$.
This applies, in particular, when the Galois group is perfect, or
is $S_d$ with $d\ge5$.
\end{corollary}
\begin{proof}
If a cubic subfield $F$ existed, its normal closure would lie in
$K$ and have a transitive cubic Galois group, necessarily $C_3$ or
$S_3$. Restriction would give the prohibited quotient. Apply
Theorem~\ref{res:thm:ambient}. Quotients of a perfect group are
perfect, whereas $C_3$ and $S_3$ are not. For $d\ge5$, the only
proper nontrivial quotient of $S_d$ is $C_2$.
\end{proof}

\begin{remark}[Composita and the remaining cubic problem]
\label{res:rem:limits}
These exclusions cannot be combined by simply taking composita.
For example, $\Q(\zeta_7)$ and $\Q(\zeta_{13})$ separately satisfy
Theorem~\ref{res:thm:galois}, but their compositum has
\[
 g_3\bigl(\Q(\zeta_{91})\bigr)
 =\frac{72}{\operatorname{lcm}(6,3)}=12.
\]
By Remark~\ref{res:rem:abelian}, that compositum contains a cyclic
cubic subfield in which $3$ splits completely.

The section excludes entire fields and all rational powers at once,
without a unit or norm hypothesis. Nevertheless, it does not eliminate
the cubic shadows already satisfying $\lambda=\pm1$, a
$\Q_3$-embedding, and Saito's real-conjugate bounds. In the cyclic
case such a shadow must have $3$ completely split. Excluding every
remaining Mills shadow would suffice for unconditional transcendence
of $A$, even without solving the stronger question for every cubic
Pisot number. No such final exclusion is obtained here.
\end{remark}


% New section: a uniform finite-prime barrier in every admissible field.
% This file uses the main manuscript's theorem environments and macros.
\section{The exact field boundary for a fixed finite set of prime divisors}
\label{res:sec:barrier}

The unit-coset result above assumes an admissible seed. Allowing a
fixed nonzero integral norm, rather than requiring norm minus one,
removes that assumption in every remaining totally real cubic field.
The resulting statement concerns divisibility by a fixed finite set
of primes; it makes no assertion that the surviving integers are prime.

For a cubic algebraic integer $\beta$, write
\[
 \lambda(\beta)=\lim_{n\to\infty}
       \Tr_{\Q(\beta)/\Q}(\beta^{3^n})\quad\hbox{in }\Z_3.
\]
Its existence and its dependence on the residual polynomial are the
three-adic convergence and classification in
Proposition~\ref{prop:fourteen} and its preceding construction.

\begin{theorem}[Every admissible field survives every fixed finite prime sieve]
\label{res:thm:equivalence}
Let $K$ be a totally real cubic field, with real embeddings
$\sigma_1,\sigma_2,\sigma_3$. The following are equivalent.
\begin{enumerate}
\item There is a field embedding $K\hookrightarrow\Q_3$.
\item For every finite set $P$ of rational primes, there are infinitely
many $\beta\in\mathcal O_K$ such that, on writing
\[
 B=\sigma_1(\beta),\qquad
 \sigma_2(\beta)=-\rho,\qquad \sigma_3(\beta)=s,
\]
one has
\begin{gather}
 B>1,\qquad 0<s<\rho<1,\qquad
 \rho<\min\{B^{-17/40},s^{17/23}\},
 \label{res:eq:strictbounds}\\
 \lambda(\beta)=1,
 \qquad
 \gcd\left(\lfloor B^m\rfloor,\prod_{p\in P}p\right)=1
 \quad\text{for every positive odd integer }m.
 \label{res:eq:avoid}
\end{gather}
\end{enumerate}
Here every $B$ in (2) is a cubic Pisot number generating $K$.
For fixed $K,P$, the elements in (2) may be chosen in one coset
$\alpha E^+(M\mathcal O_K)$, where
\[
 M=\prod_{p\in P\cup\{3\}}p,
 \qquad \Nm_{K/\Q}(\alpha)=-N,
 \qquad N\ge1,
 \qquad \gcd(N,M)=1.
\]
Thus their norm is the same integer $-N$. With
$\Delta_M$ denoting the covolume of the logarithmic lattice of
$E^+(M\mathcal O_K)$ in its first two embedding coordinates, the
number of these elements satisfying \eqref{res:eq:strictbounds} and
$B\le H$ is
\begin{equation}
\label{res:eq:count}
 \frac{3}{80\Delta_M}(\log H)^2
       +O_{K,\alpha,M}(\log H+1).
\end{equation}
The entire statement remains true with $\lambda(\beta)=-1$.
\end{theorem}

\begin{proof}
Assume first that $K\hookrightarrow\Q_3$. Then the degree-three
\'{e}tale algebra $K\otimes_\Q\Q_3$ is either
$\Q_3^3$ or $\Q_3\times E$, with $E/\Q_3$ a quadratic field.
The corresponding decomposition of $\mathcal O_K\otimes\Z_3$
allows the following prescriptions modulo three.
\begin{center}\small
\begin{tabular}{lll}
\toprule
Local algebra & Prescribed residues & Residual polynomial\\
\midrule
$\Q_3^3$ & $(1,1,-1)$ & $(X-1)^2(X+1)$\\
$\Q_3\times E$, $E$ ramified & $(-1,1)$ & $(X+1)(X-1)^2$\\
$\Q_3\times E$, $E$ unramified & $(1,i)$, $i^2=-1$ in $\F_9$ & $(X-1)(X^2+1)$\\
\bottomrule
\end{tabular}
\end{center}
In the ramified case the residue $1$ is the scalar element of
$\mathcal O_E/3\mathcal O_E$, so its multiplication operator has
characteristic polynomial $(X-1)^2$, including multiplicity. In the
unramified case $i\notin\F_3$, and its two residual conjugates
are $i$ and $-i$. Every listed element is a unit modulo three.

For each $p\in P\setminus\{3\}$ prescribe the scalar residue
$1\in\mathcal O_K/p\mathcal O_K$. The Chinese remainder theorem
gives a nonempty coset $\alpha_0+M\mathcal O_K$ satisfying all
these prescriptions. Under the three real embeddings,
$M\mathcal O_K$ is a full lattice in $\R^3$. Every translate of
such a lattice meets the open orthant $(+,-,+)$: a sufficiently
large multiple of a vector in that orthant has a lattice-coset
point within a fixed distance. Choose $\alpha$ in this intersection.
Its norm is $-N$ for a positive integer $N$, and all prescribed
residues are units, giving $\gcd(N,M)=1$.

Every element $\beta\in\alpha E^+(M\mathcal O_K)$ has the same
signs, norm, and residues as $\alpha$. Consequently, for every
positive integer $m$ and every $p\in P\setminus\{3\}$,
\[
 \Tr(\beta^m)\equiv3\pmod p.
\]
For every positive odd $m$, the three-adic prescriptions instead
give $\Tr(\beta^m)\equiv1\pmod3$: the residual roots are either
$1,1,-1$ or $1,i,-i$, with multiplicity. Their Teichm\"uller
lifts have sum one, so the classification in
Proposition~\ref{prop:fourteen} gives $\lambda(\beta)=1$ whenever
$\beta$ generates $K$.

It remains to count the elements with the required real conjugates.
The subgroup $E^+(M\mathcal O_K)$ has finite index in the unit
group. By Dirichlet's unit theorem \cite[Chapter~5]{MilneANT},
its first two logarithmic coordinates therefore form a full
rank-two lattice of covolume $\Delta_M$. Positivity removes the
torsion kernel. The corresponding coordinates of
$\alpha E^+(M\mathcal O_K)$ form a translate of this lattice.
Put $x=\log B$, $y=\log\rho$, and $L=\log N\ge0$.
The norm identity gives
\[
 \log s=L-x-y.
\]
For $a=17/40$, the inequalities
$1<B\le H$, $0<s<\rho<1$, and $\rho<B^{-a}$ are exactly
\begin{equation}
\label{res:eq:triangle}
 0<x\le\log H,\qquad \frac{L-x}{2}<y<-ax.
\end{equation}
This region has area
\[
 \frac{1-2a}{4}(\log H)^2+O_N(\log H+1)
\]
and boundary length $O_N(\log H+1)$. Counting fundamental
parallelograms, as in Theorem~\ref{transfer:thm:ray}, gives
\eqref{res:eq:count}. The leading constant is positive.

For every counted point, the dominant embedding exceeds one and
the other two have modulus less than one. Thus $B$ is Pisot and
cannot be rational; because $[K:\Q]=3$, it generates $K$.
The additional bound in \eqref{res:eq:strictbounds} follows from
the integer norm:
\[
 \rho^{40}<B^{-17}\le N^{17}B^{-17}
     =(\rho s)^{17},
 \qquad\text{so}\qquad \rho<s^{17/23}.
\]
Finally, if $m$ is positive and odd, then
$-\rho^m+s^m\in(-1,0)$. Therefore
\[
 \lfloor B^m\rfloor=\Tr_{K/\Q}(\beta^m),
\]
and the preceding residual calculations prove \eqref{res:eq:avoid}.

Conversely, a single element from (2) generates $K$ and has
$\lambda=1$. Proposition~\ref{prop:fourteen} shows that its
minimal polynomial modulo three is one of
\[
 X^2(X-1),\qquad (X-1)(X^2+1),\qquad (X-1)^2(X+1).
\]
Each has a simple linear root over $\F_3$. Hensel's lemma lifts
that root to $\Q_3$, and the resulting root of the irreducible
minimal polynomial defines an embedding $K\hookrightarrow\Q_3$.

For the variant with $\lambda=-1$, negate only the prescribed
residues at three. The real sign choice is still made independently
in the same open orthant. The odd traces are now $-1$ modulo three,
and the Teichm\"uller sum is $-1$; all other arguments are unchanged.
The reverse implication follows by negating the three residual
polynomials in the preceding paragraph.
\end{proof}

\begin{corollary}[Normalized survivors in every real interval]
\label{res:cor:normalized}
Fix an admissible field $K$, a finite prime set $P$, and a nonempty
open interval $I\subset(1,\infty)$. For every sufficiently large
integer $k$, there are elements $\beta$ from
Theorem~\ref{res:thm:equivalence} such that
\[
 B^{1/3^k}\in I,
 \qquad
 \gcd\left(\left\lfloor
    (B^{1/3^k})^{3^{k+n}}\right\rfloor,
    \prod_{p\in P}p\right)=1\quad(n\ge0).
\]
If $I$ is contained in an open interval on which a prescribed finite
list of prime floors $\lfloor t^{3^j}\rfloor$ is constant, these
normalized survivors have that same finite prime prefix.
\end{corollary}
\begin{proof}
Choose $1<a<b$ with $(a,b)\subset I$. Subtracting
\eqref{res:eq:count} at $H=b^{3^k}$ and $H=a^{3^k}$ gives
\[
 \frac{3}{80\Delta_M}3^{2k}
       \bigl((\log b)^2-(\log a)^2\bigr)+O_{K,\alpha,M,a,b}(3^k),
\]
which is positive for all sufficiently large $k$. Any resulting
$a^{3^k}<B\le b^{3^k}$ can be obtained with strict upper endpoint
as well, by choosing a smaller inner interval. The displayed tail
identity is simply
$(B^{1/3^k})^{3^{k+n}}=B^{3^n}$. The assertion about the finite
prefix follows from membership in $I$.
\end{proof}

\begin{remark}[Scope of the obstruction]
\label{res:rem:scope}
The quantifiers in Theorem~\ref{res:thm:equivalence} are
\[
 \text{for every finite }P\text{, there exists a fixed }N
 \text{ and infinitely many }\beta.
\]
The integer $N$ and the coset may depend on $P$; no one element
is asserted to avoid every prime. The normalized construction
likewise asserts neither an all-prime tail nor a prime condition
at the intervening indices between a fixed prefix and its prescribed
tail. What is ruled out is a field-uniform proof using only
Saito's strict conjugate bounds, the limiting trace condition, and
divisibility by one fixed finite set of primes. The theorem does not
rule out adaptive choices of divisors depending on $\beta$, other
arithmetic information, or an argument using the least-prime property
specific to the minimum defining Mills' constant.
\end{remark}


\section{The remaining Mills alternative}
The results yield several rigorous forms of progress, but each has
a precise quantifier. The simultaneous-window theorem closes a
residue-defined subclass. The coefficient sieve controls the size
of a possible exceptional set in boxes. The Pell certificates
exclude entire arithmetic classes of parameters in one family.
The conditional construction realizes every prescribed \emph{finite}
pattern, and the unit-lattice theorem preserves \emph{fixed finite}
local data. The new residual-field theorem supplies its own seed
in every three-adically admissible cubic field; this still preserves
a finite-divisor quantifier, not primality. None of these statements exchanges a finite quantifier
for an infinite one.

Three specific obstructions remain.
\begin{enumerate}
\item At a nonintegral three-adic limit, the prime-factor order
budget works for every fixed shift. At $\lambda=\pm1$, its key
distance from $\pm1$ is zero for the unshifted trace. These are
exactly the residual classes in which a new argument is required.
\item A coefficient set of stretched-exponentially small proportion
can still contain a particular algebraic candidate. An ambient
box estimate gives no automatic exclusion on a fixed field or on
the exponentially sparse Pell orbit. The separate Pell calculation
is needed precisely for this reason.
\item An infinite collection of algebraic numbers, each with a
long finite prime prefix, need not contain one number with an
infinite prime tower. Nor does it determine the minimum defining
Mills' constant. The minimum requires a global prime-chain argument
in addition to local polynomial information.
\end{enumerate}
Consequently the remaining full targets are unchanged: Saito's
cubic Pisot Problem~1.1, the exclusion of the Mills cubic alternative,
and irrationality of the least square-clock constant. What is
established here is a stronger collection of restricted theorems, quantitative bounds, and exact exclusions, with reusable proofs.

\paragraph{A concrete remaining target using minimality.}
For a surviving cubic candidate let $T_n=\Tr(\beta^{3^n})$.
A hypothetical algebraic least Mills constant requires both eventual
primality of $T_n$ and, for every sufficiently large $n$,
\[
 T_{n+1}=\operatorname{nextprime}(T_n^3).
\]
The latter property is an inherited observation from the supplied
manuscript [010]. To see why it holds, the prime-interval bound used
in \cite{SaitoVariants} extends every sufficiently large prime $q$
to a prime child strictly between $q^3$ and $(q+1)^3-1$.
An earlier prime child of $T_n$ could therefore be continued to an
infinite nested prime chain below the defining minimum, a contradiction.
Consequently it would suffice to prove the following restricted,
still unproved assertion: every surviving candidate either has
infinitely many composite $T_n$, or admits arbitrarily large $n$
with a prime strictly between $T_n^3$ and $T_{n+1}$.
This target uses the least-prime selection rule that the fixed-prime
survival theorem does not encode. It would settle Mills transcendence
without requiring composite floors for all cubic Pisot numbers.

\paragraph{Verification boundary.}
The central symbolic arguments were independently checked in
parallel. The residual-field proofs were checked independently,
including the ramified local algebra at three, all-rational-power
quantifiers, and the distinction between a finite-prime sieve and an
all-prime tower. Exact arithmetic replayed the two finite-field
factorizations of the balanced seed, the old 26 second-orbit Pell
witnesses and five first-orbit witnesses, the three new witnesses,
and the completeness of the new residue cover. Additional exhaustive
small-prime checks test the shifted local counts; they are consistency
checks, not substitutes for the all-degree proof. No large numerical
prime certificate from an earlier submission is needed for the new
theorems. A compact standard-library certificate is included below.

\appendix
\section{The inherited Pell reduction and the exact finite cover}
For self-contained use of the new certificate we record the earlier
reduction and cover. The substitutions $x=8u+5$, $y=4w$ transform
$4u^2+5u+23=2w^2$ into
\[
 x^2-2y^2=-343,\qquad x\equiv5\pmod8,\quad y\equiv0\pmod4.
\]
Multiplying $x+y\sqrt2$ by powers of $3-2\sqrt2$ reduces it uniquely
to $[\sqrt{343},(3+2\sqrt2)\sqrt{343})$. This bounds
$0<x<53$, $0<y<40$. Enumeration gives exactly
$(7,14),(13,16),(25,22),(35,28)$. Their four-term residue cycles
modulo eight retain only the first seed at odd powers of
$3+2\sqrt2$, and the second at even powers. These become the two
orbits from $(9,14)$ and $(1,4)$ under $T$. The latter initial point
is omitted because $u=1$ is a reducible degeneration.

For the first orbit $a_j$, the residues have periods two, ten, and
five modulo $3,19,29$ respectively:
\[
\begin{array}{c|l}
3&0,1\\
19&9,8,17,1,18,4,5,15,12,14\\
29&9,12,13,15,24.
\end{array}
\]
Even indices give trace zero modulo three. The five odd classes
$j=1,3,5,7,9\pmod{10}$ are covered by
$(p,c,n)=(29,12,41),(29,15,5),(19,4,17),(29,13,1),(19,14,7)$.
Every listed zero is beyond its root-order return budget.
For the second orbit odd $j$ give trace zero modulo three. Writing
$j=2k$, the following old rows cover all $k$ except
$0,2187\pmod{6561}$; a pair $(a,n)$ means a recurring zero at
trace index $n$ whenever $k\equiv a\pmod m$.
\begin{center}\small
\begin{tabular}{r r l}
\toprule
$p$&$m$&$(a,n)$\\\midrule
5&3&$(1,3)$\\
7&3&$(2,8)$\\
179&9&$(3,52)$\\
199&9&$(6,126)$\\
107&27&$(9,4),(18,25)$\\
487&81&$(27,366),(54,1333)$\\
4373&729&$(162,2379),(486,7323),(567,3228)$\\
20411&729&$(81,7662),(405,34139)$\\
8747&2187&$(243,7161),(324,2814),(648,105),(729,830)$\\
&&$(1053,2330),(1701,3456),(1782,4165),(2106,12465)$\\
52489&2187&$(1377,54889)$\\
39367&6561&$(3645,45679),(4374,15588)$\\
629857&2187&$(972,647491),(1458,678551)$\\
\bottomrule
\end{tabular}
\end{center}
The three new rows of Theorem~\ref{thm:pell-new} reduce the survivor
set to $k\equiv0\pmod{19683}$.

To verify the valuation formula used there, set
$(17+12\sqrt2)^{2k}=C_k+D_k\sqrt2$. Then
\[
 b_{2k}-1=\frac{13(C_k-1)+32D_k}{8},\quad
 D_{3k}=D_k(3+8D_k^2).
\]
For $3\nmid k$, reduction modulo nine gives $v_3(D_k)=1$.
Repeated cubing gives $v_3(D_k)=1+v_3(k)$ for every $k>0$.
Since $C_k^2-2D_k^2=1$ and $C_k\equiv1\pmod3$,
$v_3(C_k-1)=2v_3(D_k)$. The term involving $D_k$ therefore
determines the numerator's valuation, proving
$v_3(b_j-1)=1+v_3(j)$ for even $j>0$.

\section{A reproducible exact certificate}
The following Python 3 code uses only the standard library. It
checks primality by trial division, modular periods in the quadratic
ring, cubic trace zeros by binary polynomial powering, and the full
residue cover. It never constructs a trace of millions-digit exponent.
The algebraic reasons why these finite checks imply infinite
progressions are given in the proofs above.

\begin{lstlisting}[language=Python]
from math import isqrt

rows=[(5,3,[(1,3)]),(7,3,[(2,8)]),(179,9,[(3,52)]),(199,9,[(6,126)]),(107,27,[(9,4),(18,25)]),(487,81,[(27,366),(54,1333)]),(4373,729,[(162,2379),(486,7323),(567,3228)]),(20411,729,[(81,7662),(405,34139)]),(8747,2187,[(243,7161),(324,2814),(648,105),(729,830),(1053,2330),(1701,3456),(1782,4165),(2106,12465)]),(52489,2187,[(1377,54889)]),(39367,6561,[(3645,45679),(4374,15588)]),(629857,2187,[(972,647491),(1458,678551)])]
def v3(n):
 a=0
 while n%3==0 and n: a+=1;n//=3
 return a

def qm(x,y,p):return ((x[0]*y[0]+2*x[1]*y[1])%p,(x[0]*y[1]+x[1]*y[0])%p)
def qp(x,n,p):
 r=(1,0)
 while n:
  if n&1:r=qm(r,x,p)
  x=qm(x,x,p);n//=2
 return r

def b(j,p):
 x,y=qm((13,16),qp((17,12),j,p),p)
 return ((x-5)*pow(8,-1,p))%p

def pm(a,b,u,p):
 c=[0]*5
 for i,x in enumerate(a):
  for j,y in enumerate(b):c[i+j]=(c[i+j]+x*y)%p
 for k in (4,3):
  v=c[k];c[k]=0;c[k-1]=(c[k-1]+u*v)%p;c[k-2]=(c[k-2]+v)%p;c[k-3]=(c[k-3]-v)%p
 return c[:3]
def pp(a,n,u,p):
 r=[1,0,0]
 while n:
  if n&1:r=pm(r,a,u,p)
  a=pm(a,a,u,p);n//=2
 return r

def trace(n,u,p):
 M=(p-1)*(p*p-1)*(p**3-1); E=pow(3,n,M); A,B,C=pp([0,1,0],E,u,p)
 return (3*A+u*B+(u*u+2)*C)%p
for p,m,items in rows:
 assert p>=2 and all(p%d for d in range(2,isqrt(p)+1)),p
 assert qp((17,12),2*m,p)==(1,0),(p,m)
 for a,n in items:
  assert n>=v3((p-1)*(p*p-1)*(p**3-1)),(p,n)
  assert trace(n,b(2*a,p),p)==0,(p,m,a,n,b(2*a,p))
surv=[k for k in range(6561) if all(all(k%m!=a for a,n in items) for p,m,items in rows)]
assert surv==[0,2187],surv
for p,c,n in [(29,12,41),(29,15,5),(19,4,17),(29,13,1),(19,14,7)]:
 assert trace(n,c,p)==0
 assert n>=v3((p-1)*(p*p-1)*(p**3-1))
print('All 26 second-orbit witnesses and 5 first-orbit witnesses verified by exact arithmetic.')
print('Surviving k residues modulo6561:',surv)

# Three additional witnesses: (prime, parameter period, class, trace index).
new=[(8279983,6561,2187,4863973),
     (27398737,19683,6561,4929815),
     (15037811,19683,13122,549606)]
for p,m,a,n in new:
 assert all(p%d for d in range(2,isqrt(p)+1))
 assert qp((17,12),2*m,p)==(1,0)
 assert n>=v3((p-1)*(p*p-1)*(p**3-1))
 assert trace(n,b(2*a,p),p)==0
remaining=[k for k in range(19683)
 if all(all(k%m!=a for a,n in items) for p,m,items in rows)
 and all(k%m!=a for p,m,a,n in new)]
assert remaining==[0]
# Reduced Pell seeds and first-orbit parameter periods.
seeds=[(x,y) for x in range(1,53) for y in range(1,40)
       if x*x-2*y*y==-343]
assert seeds==[(7,14),(13,16),(25,22),(35,28)]
def orbit(seed,p,length):
 u,w=seed;out=[]
 for j in range(length):
  out.append(u%p);u,w=(17*u+12*w+10)%p,(24*u+17*w+15)%p
 assert (u,w)==(seed[0]%p,seed[1]%p)
 return out
assert orbit((9,14),3,2)==[0,1]
assert orbit((9,14),19,10)==[9,8,17,1,18,4,5,15,12,14]
assert orbit((9,14),29,5)==[9,12,13,15,24]
assert orbit((1,4),3,2)==[1,0]
print('All certificates verified; surviving k modulo19683:',remaining)

\end{lstlisting}

\section{Source map and attribution}
The supplied archive is cited as research material, not as a collection
of independently published papers. Round identifiers below refer to
\texttt{Race244\_Topic02(1).zip}; S means Showa. The earlier numbered
manuscripts refer to the previously supplied 36-paper collection.
\begin{center}\small
\begin{tabular}{p{.17\textwidth}p{.72\textwidth}}
\toprule
Source & Use and status\\\midrule
S12, S27 & Local factor blocks and uniform irreducibility, reproduced and audited.\\
S28 & Finite prime/composite patterns under H; exact multiplicities are added here.\\
S35, S36 & Complete Pell reduction and inherited finite cover; three new rows added.\\
S37 & Fixed-norm trace-zero counts and unshifted quantitative sieve; extended to all shifts and odd degrees.\\
S14 & Inherited all-norm cyclic prime-trace restrictions; used in the positive-density floor synthesis.\\
S20 & Prior norm-one trace theorem in discriminant $1524$; used only in the explicitly identified field corollary.\\
Heisei R05 & A cyclic cubic seed for finite-sieve survival; the coset-counting theorem is proved here.\\
Old [007], [008] & Previously resolved dyadic Fibonacci and nonreversible recurrence problems.\\
Old [009], [031] & Earlier local obstructions, the three-adic embedding lemma, nonreal case, qualitative density and dimension results.\\
\bottomrule
\end{tabular}
\end{center}
The other supplied topics were screened for useful connections.
The pure-cubic sixth-power observation recovers results already
subsumed by the old nonreal-cubic argument. No proved identity in
the Gieseking, integral-period, or elliptic-resolvent notes supplies
a missing prime-divisor theorem for a surviving Mills candidate.
No such speculative connection is used as a proof input.

\begin{thebibliography}{99}
\bibitem{MilneANT} J. S. Milne, \emph{Algebraic Number Theory},
version 3.08 (2020), Chapters 5--8.
\url{https://www.jmilne.org/math/CourseNotes/ANT.pdf}.
\bibitem{S14} Showa, \emph{Conductor and local restrictions on cyclic
cubic prime traces}, submitted constant-race round 14,
20 October 2026. Internal supplied research manuscript.
\bibitem{SaitoIntervals} K. Saito,
\emph{Intervals without primes near an iterated linear recurrence sequence},
International Journal of Number Theory \textbf{22} (2026), 1269--1279.
Problem numbering used here: arXiv:2504.14968v1.
\url{https://arxiv.org/html/2504.14968v1}.
\bibitem{SaitoIrrational} K. Saito, \emph{Mills' constant is irrational},
Mathematika \textbf{71} (2025), e70027.
Question numbering used here: arXiv:2404.19461v2.
\url{https://arxiv.org/html/2404.19461v2}.
\bibitem{SaitoVariants} K. Saito,
\emph{Transcendency of variants of Mills' constant},
The Ramanujan Journal \textbf{70} (2026), article 69.
Theorem numbering used here: arXiv:2508.16068v3, Theorems 1.7--1.8.
\url{https://arxiv.org/html/2508.16068v3}.
\bibitem{SaitoDimension} K. Saito,
\emph{Prime-representing functions and Hausdorff dimension},
Acta Mathematica Hungarica \textbf{165} (2021), 203--217.
Question numbering: arXiv:2102.04038v1.
\url{https://arxiv.org/abs/2102.04038}.
\bibitem{Kowalski} E. Kowalski,
\emph{The principle of the large sieve}, arXiv:math/0610021v3,
Section 5. \url{https://arxiv.org/abs/math/0610021}.
\bibitem{Milne} J. S. Milne, \emph{Algebraic Geometry},
version 6.10 (2024), B\'ezout's theorem.
\url{https://www.jmilne.org/math/CourseNotes/AG.pdf}.
\bibitem{Sutherland} A. V. Sutherland,
\emph{Global class field theory, the Chebotarev density theorem},
MIT 18.785, Lecture 28 (2021).
\url{https://math.mit.edu/classes/18.785/2021fa/LectureNotes28.pdf}.
\bibitem{Balady} S. Balady, \emph{Families of cyclic cubic fields},
arXiv:1511.04783v1 (2015).
\url{https://arxiv.org/html/1511.04783v1}.
\bibitem{Schinzel} A. Schinzel and W. Sierpi\'nski,
\emph{Sur certaines hypoth\`eses concernant les nombres premiers},
Acta Arithmetica \textbf{4} (1958), 185--208.
Hypothesis H is used as a conjectural assumption, not as a theorem.
\bibitem{Race} Showa, Heisei, and Taisho,
\emph{Constant proof race: Mills--Pisot submissions},
user-supplied research notes, 21 October 2026. In particular Showa
rounds 12, 20, 27, 28, 35, 36, 37 and Heisei round 05.
Source map above identifies the exact uses. Not external publications.
\bibitem{Keio007} Keio,
\emph{Dyadic Lucas Prime Cascades and Composite Fibonacci Shifts},
user-supplied manuscript [007], Theorem 1.1 and Corollary 1.2.
\bibitem{Keio008} Keio,
\emph{Additive Clocks, Non-Reversible Recurrences, and Transcendental
Prime-Representing Constants}, user-supplied manuscript [008],
Theorem 4.3.
\bibitem{Keio009} Keio,
\emph{Composite Terms in Sparse Cubic Pisot Floors
(Alternating Shifts, 3-Adic Limits, and Projective Return)},
user-supplied manuscript [009].
\bibitem{Keio031} Keio,
\emph{Prime-Representing Constants and Related Structures
(Encodings, Arithmetic, Geometry, and Modular Obstructions)},
user-supplied manuscript [031]. In particular Sections 14--15 and
the earlier cubic-floor and density arguments.
\end{thebibliography}
\end{document}
